\documentclass[11pt]{article}

\usepackage[margin=1in]{geometry}
\usepackage[T1]{fontenc}
\usepackage{lmodern}
\usepackage{amsmath,amssymb,amsthm}
\usepackage{mathtools}
\usepackage{enumitem}
\usepackage{microtype}
\usepackage[hidelinks]{hyperref}

\newcommand{\R}{\mathbb{R}}
\newcommand{\one}{\mathbf{1}}
\newcommand{\M}{\mathcal{M}}
\newcommand{\lexle}{\le_{\mathrm{lex}}}
\newcommand{\lexlt}{<_{\mathrm{lex}}}
\newcommand{\VB}{\mathrm{VB}}
\DeclareMathOperator{\conv}{conv}
\DeclareMathOperator{\cone}{cone}
\DeclareMathOperator{\relint}{relint}
\DeclareMathOperator{\rbd}{rbd}
\DeclareMathOperator{\aff}{aff}
\DeclareMathOperator{\supp}{supp}
\DeclareMathOperator{\ext}{ext}
\DeclareMathOperator{\cl}{cl}
\DeclareMathOperator{\inter}{int}

\theoremstyle{plain}
\newtheorem{theorem}{Theorem}[section]
\newtheorem{proposition}[theorem]{Proposition}
\newtheorem{lemma}[theorem]{Lemma}
\newtheorem{corollary}[theorem]{Corollary}
\theoremstyle{definition}
\newtheorem{definition}[theorem]{Definition}
\newtheorem{example}[theorem]{Example}
\theoremstyle{remark}
\newtheorem{remark}[theorem]{Remark}

\setlist{itemsep=2pt,topsep=4pt}

\title{Covering coherent sets of desirable gambles\\ by veto--Brier rankings}
\author{Jason Konek \and Catrin Campbell-Moore}
\date{25 October 2026}

\begin{document}
\maketitle

\begin{abstract}
Fix finitely many worlds and a full-dimensional credal set $K$. We study rankings of
models, that is, of sets of accepted gambles, evaluated world by world, on the domain of
all measurable subsets of the open cone $D_{K'}$ determined by a compact convex set $K'$
lying in the relative interior of $K$. At each world a veto--Brier ranking first charges
wrong acceptances that fall below zero under a prescribed set of probabilities, and then
minimises a Brier-type error. We show that such a ranking covers the coherent model
$D_K$ whenever the extreme points of $K$ are uniformly visible from the vertices of the
simplex: it satisfies strict guidance and independence, every model is weakly dominated
by an undominated model, every undominated model is coherent, the target is undominated,
and the domain has local freedom at the target. The proof uses compactness of the score
range, convex separation and an elementary uniqueness argument. One extreme point may be
exempt from visibility. Applications include interior $\varepsilon$-contaminations,
total-variation balls, belief functions with positive mass on the whole space and on
singletons, strictly $2$-monotone lower probabilities, smooth credal sets, every
full-dimensional credal set on two worlds, and probability intervals satisfying an
upper-bound condition. Every full-dimensional credal set is a Hausdorff limit of covered
polytopes. Finally we determine exactly which credal sets in the open simplex the
construction reaches, and identify several standard models that it does not reach.
\end{abstract}

\tableofcontents

\section{Introduction}\label{sec:intro}

A \emph{model}, in this note, is a set of gambles on finitely many worlds: the gambles it
accepts. A \emph{ranking} assigns to each world a total preorder of models, read as an
ordering by how well a model's acceptances fare if that world is actual. A dominance
argument for coherence needs a ranking with three properties: every model is weakly
dominated by an undominated model; every undominated model is coherent; and a given
coherent model, the \emph{target}, is undominated. The ranking should also meet basic
constraints. It should reward correct acceptances and penalise wrong ones (strict
guidance); it should compare models only through the gambles on which they differ
(independence); and the class of models it evaluates should leave the target free to vary
near its boundary (local freedom). When all of this holds we say that the ranking
\emph{covers} the target (Definition~\ref{def:cover}).

We give an explicit construction. Let $K$ be a full-dimensional credal set, let $D_K$ be
the set of gambles with positive expectation under every member of $K$, and let $K'$ be a
compact convex subset of the relative interior of $K$. The models are the measurable
subsets of $U=D_{K'}$, taken up to null sets. At world $i$ a \emph{veto--Brier} ranking
compares models lexicographically: first by an integral veto, which charges wrong
acceptances that fall below zero under some member of a prescribed compact set
$A_i\subseteq\conv(K'\cup\{e_i\})$, and then by the Brier-type error $-\int_X h_i\,d\mu$.

The proof follows a three-step \emph{admissibility route}.
\begin{enumerate}
  \item The vector of integrated scores ranges over a compact convex set (Lyapunov), so
  Zorn's lemma gives every model an undominated weak dominator
  (Lemma~\ref{lem:zorn}).
  \item Convex separation shows that every undominated model minimises a lexicographic
  weighting of the scores (Lemma~\ref{lem:sep}). The minimiser is unique up to null sets,
  because the relevant integrand is strictly decreasing along the all-ones direction
  (Lemma~\ref{lem:unique}); and it is coherent, because the integrand is a maximum of
  linear functions with non-negative coefficients (Lemma~\ref{lem:coh}).
  \item The target is itself one of these unique minimisers, hence undominated.
\end{enumerate}
No family of coherent models is chosen in advance: the undominated family is derived from
the ranking, and is identified exactly in Theorem~\ref{thm:family}.

The main results are the following.
\begin{itemize}
  \item Theorem~\ref{thm:cover}: if $K=\conv(K'\cup A_1\cup\dots\cup A_n)$ with
  $K'\subseteq\relint K$ and $A_i\subseteq\conv(K'\cup\{e_i\})$, then the veto--Brier
  ranking covers $K$.
  \item Theorem~\ref{thm:exempt}: one further point of the simplex, possibly on its
  boundary, may be added to the generators of $K$.
  \item Theorem~\ref{thm:vis}: the hypothesis of Theorem~\ref{thm:cover} holds exactly
  when the extreme points of $K$ are uniformly visible from the vertices of the simplex.
  We give the normal-cone form of visibility and separate strict, tangential and uniform
  visibility.
  \item Section~\ref{sec:examples}: standard classes of credal sets that are covered.
  \item Proposition~\ref{prop:dense}: every full-dimensional credal set is a Hausdorff
  limit of covered polytopes.
  \item Section~\ref{sec:limits}: for credal sets in the open simplex, the targets of
  veto--Brier rankings are exactly those with all but at most one extreme point uniformly
  visible. Pari-mutuel models, density-ratio classes and some probability intervals are
  therefore not reached by this construction.
\end{itemize}
The results of Section~\ref{sec:limits} concern this construction only. They say nothing
about rankings of other forms, or about other domains.

\section{Setting}\label{sec:setting}

\subsection{Gambles and credal sets}

Let $\Omega=\{1,\dots,n\}$ with $n\ge2$. A \emph{gamble} is a vector $h\in\R^n$; $h_i$ is
its payoff at world $i$. Write $e_1,\dots,e_n$ for the standard basis, $\one\coloneqq
(1,\dots,1)$, $\Delta\coloneqq\{m\in\R^n_+:\one\cdot m=1\}$ for the simplex, and
$\aff\Delta=\{m:\one\cdot m=1\}$. For $C\subseteq\aff\Delta$, $\relint C$ is the interior
of $C$ relative to $\aff\Delta$ and $\rbd C\coloneqq C\setminus\relint C$. A compact
convex $K\subseteq\Delta$ is \emph{full-dimensional} if $\aff K=\aff\Delta$, that is, if
$\relint K\neq\emptyset$. A \emph{credal set} is a nonempty closed convex subset of
$\Delta$. For $S\subseteq\R^n$, $\cone S$ is the set of finite non-negative combinations
of points of $S$. For $w\in\R^n$ with $\one\cdot w>0$ put $\nu(w)\coloneqq w/(\one\cdot
w)$.

For $S\subseteq\Delta$ define
\[
  D_S\coloneqq\{h\in\R^n: m\cdot h>0\text{ for all }m\in S\},\qquad
  C_S\coloneqq\{h\in\R^n: m\cdot h\ge0\text{ for all }m\in S\},
\]
with $D_\emptyset=C_\emptyset\coloneqq\R^n$. For $x\in S$ put $F^S_x\coloneqq\{h\in
C_S:x\cdot h=0\}$; when $S=K$ is fixed we write $F_x$.

\begin{lemma}\label{lem:D}
Let $S,S_1,S_2\subseteq\Delta$ be nonempty and compact.
\begin{enumerate}[label=(\alph*)]
  \item $D_S$ is an open convex cone containing every $h$ with all $h_i>0$, and
  $D_S\cap(-D_S)=\emptyset$. Moreover $D_S=D_{\conv S}$ and $D_{S_1}\cap D_{S_2}
  =D_{S_1\cup S_2}$.
  \item $\cl D_S=C_S$, and $\partial D_S=\{h:\min_{m\in S}m\cdot h=0\}$.
  \item If $S_1,S_2$ are convex and $D_{S_1}\subseteq D_{S_2}$, then $S_2\subseteq S_1$.
\end{enumerate}
\end{lemma}

\begin{proof}
(a) The function $h\mapsto\min_{m\in S}m\cdot h$ is continuous, since $|m\cdot h-m\cdot
h'|\le\max_j|h_j-h'_j|$ for $m\in\Delta$; so $D_S$ is open. Convexity, the cone property
and positivity on gambles with positive coordinates are immediate, and $m\cdot h>0$,
$m\cdot(-h)>0$ cannot both hold. A linear function has the same minimum over $S$ and
over $\conv S$, and this minimum is attained on $S$; hence $D_S=D_{\conv S}$. The last
identity is the definition.
(b) $D_S\subseteq C_S$ and $C_S$ is closed. If $h\in C_S$ then $h+t\one\in D_S$ for
$t>0$, because $m\cdot\one=1$. Hence $\cl D_S=C_S$, and $\partial D_S=C_S\setminus D_S$
since $D_S$ is open.
(c) Suppose $k\in S_2\setminus S_1$. Strict separation of $k$ from the compact convex set
$S_1$ gives $g\in\R^n$ and $c\in\R$ with $g\cdot m>c>g\cdot k$ for all $m\in S_1$. Put
$h\coloneqq g-c\one$. Since $\one\cdot m=1$ on $\Delta$, $m\cdot h>0$ on $S_1$ and
$k\cdot h<0$, so $h\in D_{S_1}\setminus D_{S_2}$.
\end{proof}

\begin{lemma}\label{lem:fulldim}
Let $K\subseteq\Delta$ be a full-dimensional compact convex set.
\begin{enumerate}[label=(\alph*)]
  \item Every affine function $\ell$ on $\aff\Delta$ has the form $\ell(m)=m\cdot h$ for
  exactly one $h\in\R^n$.
  \item $\relint K\subseteq\relint\Delta$, and $\ext K\subseteq\rbd K$.
  \item If $h\in C_K\setminus\{0\}$, then $m\cdot h>0$ for every $m\in\relint K$.
  \item If $x\in\rbd K$, then $F_x\neq\{0\}$.
\end{enumerate}
\end{lemma}

\begin{proof}
(a) If $\ell(m)=a\cdot m+c$ on $\aff\Delta$, then $\ell(m)=(a+c\one)\cdot m$ there. If
$m\cdot h=0$ for all $m\in\aff\Delta$, then $h_i=e_i\cdot h=0$ for every $i$.
(b) $\relint K$ is open in $\aff\Delta$ and contained in $\Delta$, hence in $\relint\Delta$.
A point $x\in\relint K$ is the midpoint of the segment $[x-td,x+td]\subseteq K$ for any
$d\neq0$ with $\one\cdot d=0$ and small $t>0$, so it is not extreme.
(c) If $m\mapsto m\cdot h$ is constant on $\aff\Delta$, then by (a) $h=c\one$, and $c>0$
because $h\in C_K\setminus\{0\}$. Otherwise suppose $x\in\relint K$ and $x\cdot h=0$.
Choose $d$ with $\one\cdot d=0$ and $d\cdot h\neq0$. For small $t>0$ both $x\pm td$ lie in
$K$, and one of $(x\pm td)\cdot h=\pm t\,d\cdot h$ is negative, contradicting $h\in C_K$.
(d) By the supporting hyperplane theorem in $\aff\Delta$ \cite[Thm.~11.6]{Rockafellar}
there is a non-constant affine $\ell$ on $\aff\Delta$ with $\ell\ge\ell(x)$ on $K$. By (a)
$\ell-\ell(x)$ is $m\mapsto m\cdot h$ for some $h$, and $h\neq0$. Then $h\in F_x$.
\end{proof}

We use the following standard facts of convex geometry \cite{Rockafellar}: the convex hull
of a compact subset of $\R^n$ is compact; if $S$ is compact then $\ext\conv S\subseteq S$;
every compact convex set is the convex hull of its extreme points; and if $x\in\relint K$
and $k\in K$ then $(1-s)x+sk\in\relint K$ for $0\le s<1$.

\subsection{Models and coherence}

\begin{itemize}
  \item[(M)] $\mu$ is a finite Borel measure on $\R^n$ with the same null sets as Lebesgue
  measure, and $\int|h|\,d\mu(h)<\infty$.
\end{itemize}
The standard Gaussian measure satisfies (M). Equalities and inclusions between sets of
gambles hold \emph{almost everywhere} (a.e.) if they hold up to a $\mu$-null set. In
particular every hyperplane $\{h_i=0\}$ is null. A nonempty open set is never null.

\begin{lemma}\label{lem:open}
Let $S\subseteq\Delta$ be nonempty and compact. If $O$ is open and $O\subseteq D_S$
a.e., then $O\subseteq D_S$. If $S_1,S_2\subseteq\Delta$ are nonempty, compact and
convex, and $D_{S_1}=D_{S_2}$ a.e., then $S_1=S_2$.
\end{lemma}

\begin{proof}
The open set $O\setminus\cl D_S$ is null, hence empty. So $O\subseteq\cl D_S$ and, being
open, $O\subseteq\inter\cl D_S$. For a convex set with nonempty interior the interior of
the closure is the interior \cite[Thm.~6.3]{Rockafellar}, so $O\subseteq D_S$. The second
claim follows from the first, applied to $O=D_{S_1}$ and $O=D_{S_2}$, and
Lemma~\ref{lem:D}(c).
\end{proof}

Fix a nonempty compact convex $K'\subseteq\Delta$ and put $U\coloneqq D_{K'}$. The
\emph{models} are the measurable sets $X\subseteq U$, identified when they agree a.e.;
their class is $\M_U$. A model $X$ \emph{accepts} $h$ if $h\in X$, and is otherwise
silent about $h$. At world $i$, accepting $h$ is \emph{correct} if $h_i>0$ and
\emph{wrong} if $h_i<0$. The domain $\M_U$ is decomposable: it is closed under splices
$(X\cap S)\cup(Y\setminus S)$ for measurable $S$, and under limits in measure. It is
one-sided: $U\cap(-U)=\emptyset$.

A model $X$ is \emph{coherent} if $X=D_Q$ a.e.\ for some credal set $Q$. These are the
coherent sets of desirable gambles of the form $D_Q$, taken up to null sets. By
Lemma~\ref{lem:open}, $Q$ is then unique; and since $X\subseteq U$, Lemma~\ref{lem:open}
and Lemma~\ref{lem:D}(c) give $K'\subseteq Q$.

\subsection{Rankings, dominance and covering}

A \emph{ranking} on $\M_U$ is a family $(\preceq_i)_{i\in\Omega}$ of total preorders on
$\M_U$ under which a.e.-equal models are tied. Lower is better. We write $\prec_i$ and
$\approx_i$ for the strict part and the indifference relation.

\begin{definition}\label{def:dom}
Let $X,Y\in\M_U$. $Y$ \emph{weakly dominates} $X$ if $Y\preceq_iX$ for every $i$. $Y$
\emph{dominates} $X$ if $Y$ weakly dominates $X$ and $Y\prec_iX$ for some $i$. $Y$
\emph{strictly dominates} $X$ if $Y\prec_iX$ for every $i$. $X$ is \emph{undominated} if
no model dominates it. We write $\mathcal U$ for the undominated models, $\mathcal C$
for the coherent models, and $\mathcal C^\ast\coloneqq\mathcal C\cap\mathcal U$.
\end{definition}

The axioms and conclusions are as follows.
\begin{itemize}
  \item[(G)] \emph{Strict guidance.} If $X\setminus Y\subseteq\{h_i>0\}$ a.e.,
  $Y\setminus X\subseteq\{h_i<0\}$ a.e., and $X\triangle Y$ is not null, then
  $X\prec_iY$.
  \item[(IB)] \emph{Independence.} Let $S$ be measurable and $X,Y,X',Y'\in\M_U$ with
  $X\setminus S=Y\setminus S$, $X'\setminus S=Y'\setminus S$, $X\cap S=X'\cap S$ and
  $Y\cap S=Y'\cap S$, all a.e. Then $X\preceq_iY$ if and only if $X'\preceq_iY'$.
  \item[(W)] Every model is weakly dominated by a member of $\mathcal C^\ast$.
  \item[(C)] Every model that is tied at every world with a member of $\mathcal C^\ast$
  belongs to $\mathcal C^\ast$.
  \item[(LF)] \emph{Local freedom at $D_K$.} There is an open set $N\supseteq\partial
  D_K\setminus\{0\}$ such that every measurable $X$ with $X\triangle D_K\subseteq N$ a.e.\
  belongs to $\M_U$, that is, satisfies $X\subseteq U$ a.e.
\end{itemize}

\begin{definition}\label{def:cover}
Let $K$ be a credal set. A ranking on $\M_U$ \emph{covers} $K$ if it satisfies (G) and
(IB); every model is weakly dominated by an undominated model; every undominated model is
coherent; (C) holds; $D_K\in\M_U$ and $D_K$ is undominated; and (LF) holds at $D_K$.
\end{definition}

\begin{lemma}\label{lem:dominance}
Suppose every model is weakly dominated by an undominated model and every undominated
model is coherent. Then (W) and (C) hold, $\mathcal U=\mathcal C^\ast$, and every model
outside $\mathcal C^\ast$, in particular every incoherent model, is dominated by a member
of $\mathcal C^\ast$.
\end{lemma}

\begin{proof}
(W) and $\mathcal U=\mathcal C^\ast$ are immediate. For (C), let $X\approx_iY$ for every
$i$ with $Y\in\mathcal C^\ast$. If $Z$ dominated $X$, then $Z\preceq_iX\preceq_iY$ for
every $i$ and $Z\prec_iX\preceq_iY$ for some $i$, so $Z$ would dominate $Y$. Hence
$X\in\mathcal U=\mathcal C^\ast$. Finally let $X\notin\mathcal C^\ast$ and let
$Y\in\mathcal C^\ast$ weakly dominate $X$. If $Y$ did not dominate $X$, then by
completeness $X\preceq_iY$ for every $i$, so $X$ would be tied with $Y$ at every world and
(C) would give $X\in\mathcal C^\ast$.
\end{proof}

So (C) is automatic once undominated models are coherent. The rankings below satisfy the
stronger statement that a model tied at every world with an undominated model equals it
a.e.\ (Theorem~\ref{thm:family}(d)).

\subsection{Veto--Brier rankings}

For reals, $(x_1,x_2)\lexle(y_1,y_2)$ means $x_1<y_1$, or $x_1=y_1$ and $x_2\le y_2$.

\begin{definition}\label{def:vb}
A \emph{veto profile} over $K'$ is a tuple $B=(B_1,\dots,B_n)$ of compact convex sets
with $0\in B_i\subseteq\cone(K'\cup\{e_i\})$. Put
\[
  \varphi_i(h)\coloneqq\max_{v\in B_i}(-v\cdot h),\qquad
  a_i(X)\coloneqq\int_X\varphi_i\,d\mu,\qquad
  b_i(X)\coloneqq-\int_Xh_i\,d\mu .
\]
The \emph{veto--Brier ranking} $\VB(K',B)$ on $\M_U$ is given by
\[
  X\preceq_iY\iff(a_i(X),b_i(X))\lexle(a_i(Y),b_i(Y)).
\]
The \emph{normalised veto set} of world $i$ is $\hat A_i\coloneqq\cl\{\nu(v):v\in
B_i\setminus\{0\}\}$ (empty if $B_i=\{0\}$).
\end{definition}

The error of $X$ at world $i$, weighting each wrong acceptance and each silence on a
correct gamble by $|h_i|$, is $\int_{X\cap\{h_i<0\}}|h_i|\,d\mu+\int_{(U\setminus X)
\cap\{h_i>0\}}|h_i|\,d\mu$, which equals $b_i(X)$ plus a constant. So the second level is a
Brier-type error.

The main case is the following. Let $A_1,\dots,A_n$ be compact (possibly empty) sets with
$A_i\subseteq\conv(K'\cup\{e_i\})$, and put $B_i\coloneqq\conv(\{0\}\cup A_i)$. Then
\[
  \varphi_i(h)=\max\Bigl(0,\max_{v\in A_i}(-v\cdot h)\Bigr),
\]
so the veto charges an acceptance $h$ by the amount by which its expectation falls below
zero under the least favourable member of $A_i$. We write $\VB(K',A)$ for this ranking.
Here $B_i\setminus\{0\}=\{ta:t\in(0,1],\,a\in\conv A_i\}$ and $\one\cdot a=1$, so
$\hat A_i=\conv A_i$.

\section{The admissibility route}\label{sec:route}

Throughout this section $K'\subseteq\Delta$ is nonempty, compact and convex, $U=D_{K'}$,
and $B$ is a veto profile over $K'$.

\begin{lemma}\label{lem:veto}
For every $i$:
\begin{enumerate}[label=(\alph*)]
  \item $\varphi_i$ is non-negative, convex, positively homogeneous and continuous, and
  $\varphi_i(h)\le c_i\max_j|h_j|$ with $c_i\coloneqq\max_{v\in B_i}\one\cdot v$;
  \item $\varphi_i(h)=0$ whenever $h\in U$ and $h_i\ge0$;
  \item $\varphi_i(h+t\one)\le\varphi_i(h)$ for $t\ge0$;
  \item $\hat A_i$ is a compact subset of $\conv(K'\cup\{e_i\})$, and
  $\{\varphi_i=0\}=C_{\hat A_i}$.
\end{enumerate}
\end{lemma}

\begin{proof}
Since $K'$ is convex, $\conv(K'\cup\{e_i\})=\{(1-\lambda)k+\lambda e_i:k\in K',\,
\lambda\in[0,1]\}$, and $\cone(K'\cup\{e_i\})=\{tx:t\ge0,\,x\in\conv(K'\cup\{e_i\})\}
\subseteq\R^n_+$.
(a) $\varphi_i$ is a maximum of linear functions over a compact set, and $0\in B_i$ gives
$\varphi_i\ge0$. For $v\ge0$, $|v\cdot h|\le(\one\cdot v)\max_j|h_j|$.
(b) Write $v\in B_i$ as $v=t((1-\lambda)k+\lambda e_i)$. For $h\in U$ with $h_i\ge0$,
$v\cdot h=t((1-\lambda)k\cdot h+\lambda h_i)\ge0$, so $\varphi_i(h)\le0$.
(c) $v\cdot\one\ge0$ for $v\in B_i$.
(d) Each $v\in B_i\setminus\{0\}$ is $tx$ with $t>0$ and $x\in\conv(K'\cup\{e_i\})$, and
$\nu(v)=x$. That set is closed, so it contains $\hat A_i$, which is therefore compact.
Finally $\varphi_i(h)=0$ iff $v\cdot h\ge0$ for all $v\in B_i$, iff $a\cdot h\ge0$ for all
$a\in\nu(B_i\setminus\{0\})$, iff (by continuity) $a\cdot h\ge0$ for all $a\in\hat A_i$.
\end{proof}

\begin{lemma}\label{lem:axioms}
$\VB(K',B)$ is a ranking on $\M_U$ satisfying (G) and (IB).
\end{lemma}

\begin{proof}
By Lemma~\ref{lem:veto}(a) and (M), $\varphi_i$ and $h_i$ are integrable, so $a_i,b_i$
are finite and a.e.-equal models are tied. $\lexle$ is a total order on $\R^2$, so each
$\preceq_i$ is a total preorder.

(G): Let $X,Y$ be as in (G). By Lemma~\ref{lem:veto}(b), $\varphi_i=0$ on $X\setminus Y$,
so $a_i(X)-a_i(Y)=-\int_{Y\setminus X}\varphi_i\,d\mu\le0$. Also
\[
  b_i(X)-b_i(Y)=-\int_{X\setminus Y}h_i\,d\mu+\int_{Y\setminus X}h_i\,d\mu<0,
\]
because the integrands are negative a.e.\ on their sets and $X\triangle Y$ is not null.
Hence $(a_i(X),b_i(X))\lexlt(a_i(Y),b_i(Y))$.

(IB): With $X,Y,X',Y',S$ as in (IB), $a_i(X)-a_i(Y)=\int_S(1_X-1_Y)\varphi_i\,d\mu=
a_i(X')-a_i(Y')$, and likewise for $b_i$. Since $p\lexle q$ iff $0\lexle q-p$, the two
comparisons agree.
\end{proof}

\begin{lemma}\label{lem:fubini}
Let $F:\R^n\to\R$ be continuous, and suppose $t\mapsto F(h+t\one)$ is strictly monotone
for every $h$. Then $\{F=0\}$ is null.
\end{lemma}

\begin{proof}
$\{F=0\}$ is closed. Every $h$ is uniquely $y+t\one$ with $y\in\one^\perp$ and $t\in\R$;
this linear change of coordinates multiplies Lebesgue measure by a constant. For each $y$
the section $\{t:F(y+t\one)=0\}$ has at most one point. By Fubini's theorem $\{F=0\}$ is
Lebesgue-null, hence null by (M).
\end{proof}

Write $z(X)\coloneqq(a_1(X),\dots,a_n(X),b_1(X),\dots,b_n(X))\in\R^{2n}$ and
$Z\coloneqq\{z(X):X\in\M_U\}$. For $y,y'\in\R^{2n}$ write $y\trianglelefteq y'$ if
$(y_{a_i},y_{b_i})\lexle(y'_{a_i},y'_{b_i})$ for every $i$, where $y_{a_i}$ and $y_{b_i}$
are the coordinates in the positions of $a_i$ and $b_i$. This is a partial order, because
$\lexle$ is antisymmetric. By definition, $Y$ weakly dominates $X$ iff $z(Y)\trianglelefteq
z(X)$, and $Y$ dominates $X$ iff in addition $z(Y)\neq z(X)$.

\begin{lemma}\label{lem:lyap}
$Z$ is a compact convex subset of $\R^{2n}$, and $0=z(\emptyset)\in Z$.
\end{lemma}

\begin{proof}
$z(X)=\int_Xf\,d\mu$ with $f\coloneqq(\varphi_1,\dots,\varphi_n,-h_1,\dots,-h_n)$, which
is $\mu$-integrable on $U$ by Lemma~\ref{lem:veto}(a) and (M). Each coordinate of the
vector measure $X\mapsto\int_Xf\,d\mu$ on the measurable subsets of $U$ is absolutely
continuous with respect to $\mu$, which has no atoms. By Lyapunov's convexity theorem
\cite{Lyapunov,Lindenstrauss} its range is compact and convex. Since every measurable
subset of $U$ is a model, this range is $Z$.
\end{proof}

\begin{lemma}\label{lem:zorn}
Every model is weakly dominated by an undominated model.
\end{lemma}

\begin{proof}
Fix $X_0$ and put $D\coloneqq\{y\in Z:y\trianglelefteq z(X_0)\}$. We show that every chain
$\mathcal K\subseteq D$ has a lower bound in $D$. If $\mathcal K=\emptyset$, take
$z(X_0)$. Otherwise, for $\gamma\in\mathcal K$ let $T_\gamma\coloneqq\{\gamma'\in
\mathcal K:\gamma'\trianglelefteq\gamma\}$. Since $\mathcal K$ is a chain, the intersection of finitely
many sets $T_\gamma$ contains the least of their indices, so the closed sets $\cl T_\gamma$
have the finite intersection property. As $Z$ is compact, some $y^\infty\in Z$ lies in every $\cl
T_\gamma$.

Fix $\gamma\in\mathcal K$ and a world $i$. Every $\gamma'\in T_{\gamma''}$ has
$\gamma'_{a_i}\le\gamma''_{a_i}$, so $y^\infty_{a_i}\le\gamma''_{a_i}$ for every
$\gamma''\in\mathcal K$. If $y^\infty_{a_i}<\gamma_{a_i}$, then
$(y^\infty_{a_i},y^\infty_{b_i})\lexlt(\gamma_{a_i},\gamma_{b_i})$. Otherwise
$y^\infty_{a_i}=\gamma_{a_i}$, and for $\gamma'\in T_\gamma$ we get $\gamma_{a_i}=
y^\infty_{a_i}\le\gamma'_{a_i}\le\gamma_{a_i}$. So $a_i$ is constant on $T_\gamma$, and
the lexicographic order then gives $\gamma'_{b_i}\le\gamma_{b_i}$ on $T_\gamma$. Since
$y^\infty\in\cl T_\gamma$, $y^\infty_{b_i}\le\gamma_{b_i}$. In both cases
$(y^\infty_{a_i},y^\infty_{b_i})\lexle(\gamma_{a_i},\gamma_{b_i})$.

Hence $y^\infty\trianglelefteq\gamma\trianglelefteq z(X_0)$ for all $\gamma\in\mathcal K$,
so $y^\infty\in D$ is a lower bound. By Zorn's lemma $D$ has a minimal element $y^\ast$.
Let $X^\ast$ be a model with $z(X^\ast)=y^\ast$; it weakly dominates $X_0$. If $Y$
dominated $X^\ast$, then $z(Y)\trianglelefteq y^\ast$ and $z(Y)\neq y^\ast$, so $z(Y)\in
D$ would contradict minimality.
\end{proof}

For $W_0\subseteq\Omega$ put
\[
  N(W_0)\coloneqq U\cap\bigcap_{\omega\in W_0}\{\varphi_\omega=0\}.
\]
Since $\varphi_\omega\ge0$, a model $X$ has $a_\omega(X)=0$ iff $X\subseteq
\{\varphi_\omega=0\}$ a.e. A triple $(W_0,\beta,\alpha)$ is \emph{admissible} if
$W_0\subseteq\Omega$, $\beta\in\R^n_+\setminus\{0\}$ with $\supp\beta\subseteq W_0$, and
$\alpha\in\R_+^{\Omega\setminus W_0}$. For such a triple put
\[
  J(X)\coloneqq\sum_{\omega\in W_0}\beta_\omega b_\omega(X)
    +\sum_{\omega\notin W_0}\alpha_\omega a_\omega(X),\qquad
  g(h)\coloneqq\sum_{\omega\notin W_0}\alpha_\omega\varphi_\omega(h)-\beta\cdot h .
\]
Then $J(X)=\int_Xg\,d\mu$ whenever $X\subseteq N(W_0)$ a.e.

\begin{lemma}\label{lem:sep}
Let $X^\ast$ be undominated and $W_0\coloneqq\{\omega:a_\omega(X^\ast)=0\}$. Then $W_0\neq
\emptyset$, $X^\ast\subseteq N(W_0)$ a.e., and there are $\beta,\alpha$ with
$(W_0,\beta,\alpha)$ admissible such that $X^\ast$ minimises $J$ over the models $X\subseteq
N(W_0)$.
\end{lemma}

\begin{proof}
If $W_0=\emptyset$, every $a_\omega(X^\ast)>0=a_\omega(\emptyset)$, so $\emptyset$ strictly
dominates $X^\ast$. Hence $W_0\neq\emptyset$, and $X^\ast\subseteq N(W_0)$ a.e.\ by the
remark above.

Let $Z'\coloneqq\{y\in Z:y_{a_\omega}=0\text{ for }\omega\in W_0\}$, a convex set
containing $z(X)$ for every model $X\subseteq N(W_0)$, and $y^\ast\coloneqq z(X^\ast)\in
Z'$. Let $\mathrm{pr}$ be the projection of $\R^{2n}$ onto the coordinates $y_{b_\omega}$
($\omega\in W_0$) and $y_{a_\omega},y_{b_\omega}$ ($\omega\notin W_0$), and let
\[
  P\coloneqq\{d:d_{b_\omega}<0\text{ for }\omega\in W_0,\ d_{a_\omega}<0\text{ for }
  \omega\notin W_0\}
\]
in the target of $\mathrm{pr}$, a nonempty open convex cone. If some $y\in Z'$ had
$\mathrm{pr}(y-y^\ast)\in P$, the model realising $y$ would be strictly better than $X^\ast$ at
every world: at $\omega\in W_0$ its veto is equal (zero) and its error strictly smaller,
and at $\omega\notin W_0$ its veto is strictly smaller. So the convex set $E\coloneqq
\mathrm{pr}(Z'-y^\ast)$ is disjoint from $P$.

By the separation theorem \cite[Thm.~11.3]{Rockafellar} there is a nonzero linear
functional $L$ with $\sup_PL\le\inf_EL$. Since $0\in E$ and $P$ is a cone with $0\in\cl P$,
$\sup_PL=0$; so $L\le0$ on $\cl P$ and $L\ge0$ on $E$. Write
\[
  L(d)=\sum_{\omega\in W_0}\beta_\omega d_{b_\omega}+\sum_{\omega\notin W_0}
  (\alpha_\omega d_{a_\omega}+\gamma_\omega d_{b_\omega}).
\]
Testing $L\le0$ on the vectors of $\cl P$ with one coordinate equal to $-1$, and on
$\pm$ the unit vectors in the coordinates $d_{b_\omega}$ ($\omega\notin W_0$), which are
unconstrained in $\cl P$, gives $\beta\ge0$, $\alpha\ge0$ and $\gamma=0$. Extend $\beta$
by zero outside $W_0$. Then $L(\mathrm{pr}(y-y^\ast))=J(y)-J(y^\ast)\ge0$ for $y\in Z'$, where
$J$ is read as a function of the score vector. So $X^\ast$ minimises $J$ over the models
$X\subseteq N(W_0)$.

It remains to show $\beta\neq0$. Otherwise $\alpha\neq0$ and $J=\sum_{\omega\notin W_0}
\alpha_\omega a_\omega\ge0$ on $Z'$, with $J(0)=0$ and $0=z(\emptyset)\in Z'$. Hence
$J(y^\ast)=0$, so $\alpha_\omega a_\omega(X^\ast)=0$ for every $\omega\notin W_0$. But
$a_\omega(X^\ast)>0$ for $\omega\notin W_0$, so $\alpha=0$, a contradiction.
\end{proof}

\begin{lemma}\label{lem:unique}
Let $(W_0,\beta,\alpha)$ be admissible. Then $\{g=0\}$ is null, and the models $X\subseteq
N(W_0)$ minimising $\int_Xg\,d\mu$ are exactly those equal a.e.\ to
\[
  X^\ast(W_0,\beta,\alpha)\coloneqq N(W_0)\cap\{g<0\}.
\]
\end{lemma}

\begin{proof}
$g$ is continuous and integrable by Lemma~\ref{lem:veto}(a). By Lemma~\ref{lem:veto}(c),
for $t>0$,
\[
  g(h+t\one)\le g(h)-t\,\one\cdot\beta<g(h),
\]
since $\one\cdot\beta>0$. Lemma~\ref{lem:fubini} makes $\{g=0\}$ null. For $X\subseteq
N\coloneqq N(W_0)$,
\[
  \int_Xg\,d\mu=\int_{X\cap\{g<0\}}g\,d\mu+\int_{X\cap\{g>0\}}g\,d\mu\ \ge\
  \int_{N\cap\{g<0\}}g\,d\mu,
\]
with equality iff $(N\cap\{g<0\})\setminus X$ and $X\cap\{g>0\}$ are null, that is, iff
$X=N\cap\{g<0\}$ a.e.
\end{proof}

\begin{lemma}\label{lem:coh}
Let $(W_0,\beta,\alpha)$ be admissible, and put
\[
  W\coloneqq\Bigl\{\beta+\sum_{\omega\notin W_0}\alpha_\omega v_\omega:v_\omega\in
  B_\omega\Bigr\},\qquad\hat W\coloneqq\nu(W).
\]
Then $\hat W\subseteq\Delta$ is compact, and $X^\ast(W_0,\beta,\alpha)=D_{K(W_0,\beta,
\alpha)}$ a.e., where
\[
  K(W_0,\beta,\alpha)\coloneqq\conv\Bigl(K'\cup\bigcup_{\omega\in W_0}\hat A_\omega\cup
  \hat W\Bigr)
\]
is a credal set containing $K'$. In particular $X^\ast(W_0,\beta,\alpha)$ is coherent.
\end{lemma}

\begin{proof}
Since $\alpha\ge0$, a sum of maxima over independent choices is the maximum of the sums:
\[
  g(h)=\max_{w\in W}(-w\cdot h).
\]
$W$ is compact, as a continuous image of a product of compact sets. Each $w\in W$ is
non-negative with $\one\cdot w\ge\one\cdot\beta>0$, so $\nu$ is continuous on $W$ and
$\hat W\subseteq\Delta$ is compact. As the maximum is attained, $\{g<0\}=\{h:w\cdot h>0
\text{ for all }w\in W\}=D_{\hat W}$.

Let $\omega\in W_0$. By Lemma~\ref{lem:veto}(d), $\{\varphi_\omega=0\}=C_{\hat
A_\omega}$. If $\hat A_\omega=\emptyset$, this is $\R^n=D_{\hat A_\omega}$. Otherwise
$\psi_\omega(h)\coloneqq\min_{a\in\hat A_\omega}a\cdot h$ is continuous with
$\psi_\omega(h+t\one)=\psi_\omega(h)+t$, so $\{\psi_\omega=0\}$ is null by
Lemma~\ref{lem:fubini}, and $C_{\hat A_\omega}=D_{\hat A_\omega}\cup\{\psi_\omega=0\}$.
Either way $\{\varphi_\omega=0\}=D_{\hat A_\omega}$ a.e. By Lemma~\ref{lem:D}(a),
\[
  X^\ast(W_0,\beta,\alpha)=D_{K'}\cap\bigcap_{\omega\in W_0}D_{\hat A_\omega}\cap
  D_{\hat W}=D_{K'\cup\bigcup_{\omega\in W_0}\hat A_\omega\cup\hat W}=
  D_{K(W_0,\beta,\alpha)}\quad\text{a.e.}
\]
The union is a finite union of compact subsets of $\Delta$, so its convex hull is a
credal set.
\end{proof}

\begin{theorem}\label{thm:family}
Let $B$ be a veto profile over $K'$. The ranking $\VB(K',B)$ on $\M_U$ has the following
properties.
\begin{enumerate}[label=(\alph*)]
  \item It satisfies (G) and (IB).
  \item Every model is weakly dominated by an undominated model.
  \item A model is undominated if and only if it equals $X^\ast(W_0,\beta,\alpha)$ a.e.\
  for some admissible triple. Every undominated model is therefore coherent, and equals
  $D_{K(W_0,\beta,\alpha)}$ a.e.
  \item If a model is tied at every world with an undominated model $X$, it equals $X$
  a.e.
\end{enumerate}
Consequently (W) and (C) hold, $\mathcal U=\mathcal C^\ast$, and every incoherent model is
dominated by a coherent undominated model.
\end{theorem}

\begin{proof}
(a) is Lemma~\ref{lem:axioms} and (b) is Lemma~\ref{lem:zorn}.

(c) Let $X^\ast$ be undominated. Lemma~\ref{lem:sep} gives an admissible triple such that
$X^\ast\subseteq N(W_0)$ minimises $\int_Xg\,d\mu$ over the models $X\subseteq N(W_0)$. By
Lemma~\ref{lem:unique}, $X^\ast=X^\ast(W_0,\beta,\alpha)$ a.e., and
Lemma~\ref{lem:coh} identifies it.

Conversely, let $(W_0,\beta,\alpha)$ be admissible, $X^\ast\coloneqq X^\ast(W_0,\beta,
\alpha)$, and let $Y$ weakly dominate $X^\ast$. For $\omega\in W_0$, $a_\omega(X^\ast)=0$,
so $a_\omega(Y)\le0$; hence $a_\omega(Y)=0$, and the lexicographic order then gives
$b_\omega(Y)\le b_\omega(X^\ast)$. For $\omega\notin W_0$, $a_\omega(Y)\le
a_\omega(X^\ast)$. Thus $Y\subseteq N(W_0)$ a.e.\ and, as the weights are non-negative,
$J(Y)\le J(X^\ast)$. By Lemma~\ref{lem:unique}, $Y=X^\ast$ a.e. So every weak dominator of
$X^\ast$ equals it a.e., and $X^\ast$ is undominated.

(d) Let $Y$ be tied with $X$ at every world. Then $z(Y)=z(X)$. By (c),
$X=X^\ast(W_0,\beta,\alpha)$ a.e.\ for an admissible triple. Then $a_\omega(Y)=0$ for
$\omega\in W_0$, so $Y\subseteq N(W_0)$ a.e., and $J(Y)=J(X)$. By
Lemma~\ref{lem:unique}, $Y=X$ a.e.

The final sentence follows from (b), (c) and Lemma~\ref{lem:dominance}.
\end{proof}

\section{The covering theorem}\label{sec:cover}

\begin{lemma}\label{lem:LF}
Let $K\subseteq\Delta$ be full-dimensional, compact and convex, and let $K'\subseteq\Delta$
be nonempty, compact and convex, with $U=D_{K'}$.
\begin{enumerate}[label=(\alph*)]
  \item If $K'\subseteq\relint K$, then $D_K\subseteq U$ and $\partial D_K\setminus\{0\}
  \subseteq U$. Hence $D_K\in\M_U$ and (LF) holds at $D_K$ with $N=U$.
  \item Conversely, if $D_K\subseteq U$ a.e.\ and (LF) holds at $D_K$, then $K'\subseteq
  \relint K$.
\end{enumerate}
\end{lemma}

\begin{proof}
(a) $K'\subseteq K$ gives $D_K\subseteq D_{K'}$. Let $h\in\partial D_K\setminus\{0\}$. By
Lemma~\ref{lem:D}(b), $h\in C_K\setminus\{0\}$, so Lemma~\ref{lem:fulldim}(c) gives
$k'\cdot h>0$ for every $k'\in K'$, that is, $h\in U$. $U$ is open, and every measurable
$X$ with $X\triangle D_K\subseteq U$ a.e.\ satisfies $X\subseteq U$ a.e.

(b) By Lemma~\ref{lem:open}, $D_K\subseteq U$, and Lemma~\ref{lem:D}(c) gives
$K'\subseteq K$. Let $N$ be as in (LF). By (LF) the set $D_K\cup N$ belongs to $\M_U$, so
$N\subseteq U$ a.e., and $N\subseteq U$ by Lemma~\ref{lem:open}. Suppose $k'\in K'\cap\rbd K$. By
Lemma~\ref{lem:fulldim}(d) there is $h\in F_{k'}\setminus\{0\}$. Then $h\in C_K\setminus
D_K=\partial D_K$, so $h\in N\subseteq U$; but $k'\cdot h=0$, so $h\notin U$.
\end{proof}

\begin{theorem}\label{thm:cover}
Let $K\subseteq\Delta$ be a full-dimensional credal set. Suppose that $K'\subseteq\relint
K$ is nonempty, compact and convex, that $A_1,\dots,A_n$ are compact (possibly empty)
with $A_i\subseteq\conv(K'\cup\{e_i\})$, and that
\[
  K=\conv(K'\cup A_1\cup\dots\cup A_n).
\]
Then $\VB(K',A)$ covers $K$ on $\M_{D_{K'}}$. Its undominated models are exactly the
coherent models $D_{K(W_0,\beta,\alpha)}$ of Theorem~\ref{thm:family}, and a model tied at
every world with an undominated model equals it a.e.
\end{theorem}

\begin{proof}
Theorem~\ref{thm:family} gives (G), (IB), the existence of undominated weak dominators,
coherence of undominated models, (C), and the last two claims. Lemma~\ref{lem:LF}(a)
gives $D_K\in\M_U$ and (LF).

\emph{Target protection.} Fix $p\in K'$. The triple $(\Omega,p,\emptyset)$ is admissible
(there are no $\alpha$-weights), and $W=\{p\}$, $\hat W=\{p\}$. Since $\hat
A_\omega=\conv A_\omega$ (Section~\ref{sec:setting}),
\[
  K(\Omega,p,\emptyset)=\conv\Bigl(K'\cup\bigcup_\omega\conv A_\omega\cup\{p\}\Bigr)
  =\conv(K'\cup A_1\cup\dots\cup A_n)=K .
\]
By Lemma~\ref{lem:coh} and Theorem~\ref{thm:family}(c), $D_K=X^\ast(\Omega,p,\emptyset)$
a.e.\ is undominated.
\end{proof}

\begin{remark}\label{rem:direct}
The target step can also be seen directly. If $Y$ weakly dominates $D_K$, then
$a_i(Y)\le a_i(D_K)=0$ for every $i$, so $Y\subseteq N(\Omega)=D_K$ a.e.\ by the
computation above. With the vetoes tied, $b_i(Y)\le b_i(D_K)$ for every $i$. For $p\in K$,
$\sum_ip_i(b_i(Y)-b_i(D_K))=\int_{D_K\setminus Y}p\cdot h\,d\mu$, which is positive unless
$D_K\setminus Y$ is null. Hence $Y=D_K$ a.e.
\end{remark}

\begin{corollary}\label{cor:contam}
Let $p\in\relint\Delta$, $0<\varepsilon\le1$ and $K\coloneqq(1-\varepsilon)p+
\varepsilon\Delta$, with vertices $v_i\coloneqq(1-\varepsilon)p+\varepsilon e_i$. For
$0\le s<1$ let $K'\coloneqq p+s(K-p)$ and $A_i\coloneqq\{v_i\}$. Then $\VB(K',A)$ covers
$K$ on $\M_{D_{K'}}$, in every number $n\ge2$ of worlds. At world $i$ the ranking compares
\[
  \Bigl(\int_X\bigl[-(1-\varepsilon)\,p\cdot h-\varepsilon h_i\bigr]_+d\mu(h),\
  -\int_Xh_i\,d\mu(h)\Bigr)
\]
lexicographically. The case $\varepsilon=1$ is the vacuous model $K=\Delta$.
\end{corollary}

\begin{proof}
The map $m\mapsto(1-\varepsilon)p+\varepsilon m$ is an affine bijection of $\aff\Delta$
mapping $\Delta$ onto $K$ and fixing $p$, so $K$ is full-dimensional and $p\in\relint K$.
Hence $K'\subseteq\relint K$. Also $v_i\in\conv(\{p\}\cup\{e_i\})\subseteq\conv(K'\cup
\{e_i\})$, and $K=\conv\{v_1,\dots,v_n\}=\conv(K'\cup A_1\cup\dots\cup A_n)$ because
$K'\subseteq K$. Apply Theorem~\ref{thm:cover}; the veto is $\max(0,-v_i\cdot h)$.
\end{proof}

\section{One exempt point}\label{sec:exempt}

\begin{theorem}\label{thm:exempt}
Let $K'\subseteq\Delta$ be nonempty, compact and convex, let $A_1,\dots,A_n$ be compact
with $A_i\subseteq\conv(K'\cup\{e_i\})$, and let $p\in\Delta$ be arbitrary. Put
\[
  K\coloneqq\conv(K'\cup A_1\cup\dots\cup A_n\cup\{p\}).
\]
If $K$ is full-dimensional and $K'\subseteq\relint K$, then $\VB(K',A)$ covers $K$ on
$\M_{D_{K'}}$.
\end{theorem}

\begin{proof}
The triple $(\Omega,p,\emptyset)$ is admissible, since $p\in\R^n_+\setminus\{0\}$, and
$\hat W=\{p\}$. As in the proof of Theorem~\ref{thm:cover}, $K(\Omega,p,\emptyset)=K$, so
$D_K$ is undominated by Theorem~\ref{thm:family}(c). Lemma~\ref{lem:LF}(a) gives $D_K\in
\M_U$ and (LF), and Theorem~\ref{thm:family} gives the remaining requirements.
\end{proof}

\begin{remark}\label{rem:exempt}
The exempt point $p$ may lie on the boundary of $\Delta$. The set
$K_0\coloneqq\conv(K'\cup A_1\cup\dots\cup A_n)$ need not be full-dimensional, and $K'$
need not lie in $\relint K_0$. The ranking does not depend on $p$: a single ranking
$\VB(K',A)$ covers every $\conv(K_0\cup\{p\})$, $p\in\Delta$, that is full-dimensional and
contains $K'$ in its relative interior.
\end{remark}

\section{Visibility}\label{sec:vis}

Throughout this section $K\subseteq\Delta$ is a full-dimensional credal set.

\begin{definition}\label{def:vis}
Let $x\in K$ and $j\in\Omega$.
\begin{enumerate}[label=(\alph*)]
  \item $x$ is \emph{visible from $j$} if $x=e_j$, or the open ray $\{x+s(x-e_j):s>0\}$
  meets $\relint K$. $x$ is \emph{visible} if it is visible from some world.
  \item $x$ is \emph{weakly visible from $j$} if $h_j\le0$ for every $h\in F_x$, and
  \emph{tangential from $j$} if it is weakly visible but not visible from $j$.
  \item For $\delta>0$ let $K_\delta\coloneqq\{k\in K:\bar B(k,\delta)\cap\aff\Delta
  \subseteq K\}$, where $\bar B(k,\delta)$ is the closed Euclidean ball. A set $S\subseteq
  K$ is \emph{uniformly visible} if, for some $\delta>0$,
  \[
    S\subseteq\bigcup_{j\in\Omega}\conv(K_\delta\cup\{e_j\}).
  \]
\end{enumerate}
\end{definition}

\begin{lemma}\label{lem:visbasic}
\begin{enumerate}[label=(\alph*)]
  \item Let $x\in K$ and $x\neq e_j$. Then $x$ is visible from $j$ if and only if
  $x=(1-t)k+te_j$ for some $k\in\relint K$ and $t\in[0,1)$.
  \item Each $K_\delta$ is compact, convex and contained in $\relint K$; $K_\delta
  \supseteq K_{\delta'}$ for $\delta\le\delta'$; $K_\delta\neq\emptyset$ for small
  $\delta$; and every compact subset of $\relint K$ lies in some $K_\delta$.
  \item Let $x=(1-t)k+te_j$ with $k\in K_\delta$ and $t\in[0,1)$. Then every $y\in K$ with
  $|y-x|<(1-t)\delta/2$ lies in $\conv(K_{\delta/2}\cup\{e_j\})$.
\end{enumerate}
\end{lemma}

\begin{proof}
(a) If $k\coloneqq x+s(x-e_j)\in\relint K$ with $s>0$, then $x=\frac1{1+s}k+\frac
s{1+s}e_j$. Conversely let $x=(1-t)k+te_j$ with $k\in\relint K$. If $t>0$, then
$k=x+\frac t{1-t}(x-e_j)$ lies on the open ray. If $t=0$, then $x\in\relint K$, which is
open in $\aff\Delta$ and contains the points of the ray near $x$.

(b) Let $d(k)\coloneqq\inf\{|k-y|:y\in\aff\Delta\setminus K\}$, a $1$-Lipschitz function.
Since $K$ is closed, $K_\delta=\{k\in K:d(k)\ge\delta\}$, which is compact. If $k,k'\in
K_\delta$, $\lambda\in[0,1]$ and $y=\lambda k+(1-\lambda)k'+u$ with $|u|\le\delta$ and
$\one\cdot u=0$, then $y=\lambda(k+u)+(1-\lambda)(k'+u)\in K$; so $K_\delta$ is convex.
Clearly $K_\delta\subseteq\relint K$ and the sets increase as $\delta$ decreases. $d$ is
positive and continuous on the nonempty set $\relint K$, so it has a positive minimum on
any compact subset of $\relint K$.

(c) Put $k_y\coloneqq(y-te_j)/(1-t)\in\aff\Delta$. Then $|k_y-k|=|y-x|/(1-t)<\delta/2$, so
$\bar B(k_y,\delta/2)\subseteq\bar B(k,\delta)$ and $k_y\in K_{\delta/2}$. Moreover
$y=(1-t)k_y+te_j$.
\end{proof}

\begin{lemma}\label{lem:V1}
Let $x\in K$ and $x\neq e_j$. Then $x$ is visible from $j$ if and only if $h_j<0$ for
every nonzero $h\in F_x$. In particular, every point of $\relint K$ is visible from every
world, and visible points are weakly visible.
\end{lemma}

\begin{proof}
($\Rightarrow$) Let $k=x+s(x-e_j)\in\relint K$ with $s>0$, and $h\in F_x\setminus\{0\}$.
Then $k\cdot h=(1+s)\,x\cdot h-s\,h_j=-s\,h_j$, and $k\cdot h>0$ by
Lemma~\ref{lem:fulldim}(c). So $h_j<0$.

($\Leftarrow$) Suppose the open ray $R$ misses $\relint K$. Both are convex subsets of
$\aff\Delta$ and $\relint K$ is relatively open and nonempty. By separation in $\aff\Delta$
\cite[Thm.~11.3]{Rockafellar} there are a non-constant affine $\ell$ on $\aff\Delta$ and
$c\in\R$ with $\ell\ge c$ on $\relint K$, hence on $K$, and $\ell\le c$ on $R$. Letting
$s\to0$ along $R$ gives $\ell(x)\le c$, and $x\in K$ gives $\ell(x)\ge c$; so $\ell(x)=c$.
By Lemma~\ref{lem:fulldim}(a), $\ell-c$ is $m\mapsto m\cdot h$ for some $h\neq0$. Then
$h\in F_x$, and on $R$, $(x+s(x-e_j))\cdot h=-s\,h_j\le0$, so $h_j\ge0$.

For $x\in\relint K$, $F_x=\{0\}$ by Lemma~\ref{lem:fulldim}(c).
\end{proof}

\begin{theorem}\label{thm:vis}
The following are equivalent.
\begin{enumerate}[label=(\roman*)]
  \item There are $K'$ and $A_1,\dots,A_n$ satisfying the hypotheses of
  Theorem~\ref{thm:cover}.
  \item $\ext K$ is uniformly visible.
  \item Both of the following hold:
  \begin{enumerate}[label=(\alph*)]
    \item every point of $\cl(\ext K)$ that is not a vertex of $\Delta$ is visible;
    \item no vertex $e_i$ of $\Delta$ that lies in $K$ is an accumulation point of
    $\ext K$.
  \end{enumerate}
\end{enumerate}
Under these conditions $K$ is covered by $\VB(K_\delta,A)$ with
$A_j\coloneqq K\cap\conv(K_\delta\cup\{e_j\})$, for $\delta$ as in (ii) and small enough
that $K_\delta\neq\emptyset$.
\end{theorem}

\begin{proof}
(i)$\Rightarrow$(ii). $\ext K\subseteq K'\cup\bigcup_iA_i$, since $K$ is the convex hull of
this compact set, and $K'\cap\ext K=\emptyset$ by Lemma~\ref{lem:fulldim}(b). By
Lemma~\ref{lem:visbasic}(b), $K'\subseteq K_\delta$ for some $\delta>0$, so $\ext
K\subseteq\bigcup_i\conv(K_\delta\cup\{e_i\})$.

(ii)$\Rightarrow$(i). Decreasing $\delta$ preserves (ii), so we may assume $K_\delta\neq
\emptyset$. Put $K'\coloneqq K_\delta\subseteq\relint K$ and $A_j\coloneqq K\cap
\conv(K_\delta\cup\{e_j\})$, which is compact. Then $\ext K\subseteq\bigcup_jA_j$, and
$K=\conv\ext K\subseteq\conv(K'\cup\bigcup_jA_j)\subseteq K$. This also proves the last
sentence of the theorem, by Theorem~\ref{thm:cover}.

(ii)$\Rightarrow$(iii). Let $x\in\cl(\ext K)$, $x=\lim v_m$ with $v_m\in\ext K$, and write
$v_m=(1-t_m)k_m+t_me_{j_m}$ with $k_m\in K_\delta$ and $t_m\in[0,1]$. Passing to a
subsequence, $j_m=j$, $t_m\to t$ and $k_m\to k\in K_\delta$, so $x=(1-t)k+te_j$. For (a):
if $x$ is not a vertex of $\Delta$, then $t<1$, and $x$ is visible from $j$ by
Lemma~\ref{lem:visbasic}(a). For (b): suppose $e_i\in K$ and $v_m\to e_i$ with $v_m\neq
e_i$. Then $e_i=(1-t)k+te_j$ with $k\in\relint K\subseteq\relint\Delta$
(Lemma~\ref{lem:fulldim}(b)); since $k$ has full support and $n\ge2$, this forces $t=1$
and $j=i$. So $v_m=(1-t_m)k_m+t_me_i$ with $t_m<1$ (as $v_m\ne e_i$) and $t_m>0$ (as
$v_m\notin\relint K$). Then $v_m$ lies in the open segment between $k_m\in K$ and $e_i\in
K$, so it is not extreme, a contradiction.

(iii)$\Rightarrow$(ii). For $x\in\cl(\ext K)$ not a vertex of $\Delta$, (a) and
Lemma~\ref{lem:visbasic}(a) give $x=(1-t)k+te_j$ with $k\in\relint K$, $t\in[0,1)$; let
$\delta(x)\coloneqq d(k)>0$, so $k\in K_{\delta(x)}$. By Lemma~\ref{lem:visbasic}(c) the open
ball $O_x$ of radius $(1-t)\delta(x)/2$ about $x$ satisfies $O_x\cap K\subseteq
\conv(K_{\delta(x)/2}\cup\{e_j\})$. For a vertex $e_i\in\cl(\ext K)$, (b) gives an open
ball $O_{e_i}$ about $e_i$ meeting $\ext K$ only in $e_i$, and $e_i\in\conv(K_\delta\cup
\{e_i\})$ for every $\delta$ with $K_\delta\ne\emptyset$. The compact set $\cl(\ext K)$ is
covered by finitely many of these balls. Let $\delta$ be the least of the corresponding
numbers $\delta(x)/2$, decreased if necessary so that $K_\delta\neq\emptyset$. Since
$K_{\delta'}\subseteq K_\delta$ for $\delta'\ge\delta$, every extreme point lies in some
$\conv(K_\delta\cup\{e_j\})$.
\end{proof}

\begin{corollary}\label{cor:exemptvis}
If, for some $p\in K$, the set $\ext K\setminus\{p\}$ is uniformly visible, then $K$ is
covered by $\VB(K_\delta,A)$ with $A_j\coloneqq K\cap\conv(K_\delta\cup\{e_j\})$, for
suitable $\delta>0$.
\end{corollary}

\begin{proof}
Choose $\delta$ for $\ext K\setminus\{p\}$ with $K_\delta\neq\emptyset$, and put
$K'\coloneqq K_\delta$. Then $\ext K\subseteq\bigcup_jA_j\cup\{p\}$, so
$K=\conv(K'\cup\bigcup_jA_j\cup\{p\})$, and $K'\subseteq\relint K$. Apply
Theorem~\ref{thm:exempt}.
\end{proof}

For polytopes, visibility can be read off the facet normals.

\begin{lemma}\label{lem:poly}
Suppose $K=\{m\in\aff\Delta:g_f\cdot m\ge0\text{ for }f\in I\}$ with $I$ finite. Then
$C_K=\cone\{g_f:f\in I\}+\R_+\one$, and for $x\in K$,
\[
  F_x=\cone\{g_f:f\in I,\ g_f\cdot x=0\}.
\]
Hence $x\neq e_j$ is visible from $j$ if and only if $(g_f)_j<0$ for every $f$ with
$g_f\neq0$ and $g_f\cdot x=0$.
\end{lemma}

\begin{proof}
By the affine Farkas lemma \cite[Cor.~7.1h]{Schrijver}, applied to the system $g_f\cdot
m\ge0$, $\one\cdot m\ge1$, $-\one\cdot m\ge-1$, the inequality $h\cdot m\ge0$ holds on
$K$ iff $h=\sum_f\lambda_fg_f+(\eta_+-\eta_-)\one$ with $\lambda,\eta_\pm\ge0$ and
$\eta_+-\eta_-\ge0$. This is the formula for $C_K$. If such an $h$ has $x\cdot h=0$, then
$\sum_f\lambda_f\,g_f\cdot x+(\eta_+-\eta_-)=0$ with non-negative terms, so $\lambda_f=0$
unless $g_f\cdot x=0$, and the $\one$-term vanishes. The visibility criterion follows from
Lemma~\ref{lem:V1}: a nonzero element of $F_x$ has a positive coefficient on some nonzero
active $g_f$, and each nonzero active $g_f$ lies in $F_x$.
\end{proof}

Every polytope $K\subseteq\Delta$ has such a description: an affine inequality $a\cdot
m\ge c$ on $\aff\Delta$ is $(a-c\one)\cdot m\ge0$.

\begin{corollary}\label{cor:poly}
Let $K$ be a full-dimensional polytope. If every vertex of $K$, with at most one
exception, is a vertex of $\Delta$ or visible, then $K$ is covered by a veto--Brier
ranking. If there is no exception, the hypotheses of Theorem~\ref{thm:cover} hold.
\end{corollary}

\begin{proof}
A visible vertex $v\neq e_j$ has the form $(1-t)k+te_j$ with $k\in\relint K$, so
$v\in\conv(K_{d(k)}\cup\{e_j\})$; and $e_j\in\conv(K_\delta\cup\{e_j\})$. Taking the least
of finitely many depths, $\ext K$ minus the exception is uniformly visible. Apply
Corollary~\ref{cor:exemptvis} or Theorem~\ref{thm:vis}.
\end{proof}

\begin{corollary}\label{cor:smooth}
Let $K\subseteq\relint\Delta$ be full-dimensional, and suppose $K$ has exactly one
supporting hyperplane (in $\aff\Delta$) at each point of $\rbd K$. Then every point of $K$
is visible, the hypotheses of Theorem~\ref{thm:cover} hold, and $K$ is covered.
\end{corollary}

\begin{proof}
Let $x\in\rbd K$. By Lemma~\ref{lem:fulldim}(d) there is $h_x\in F_x\setminus\{0\}$. If
$h\in F_x\setminus\{0\}$, the affine functions $m\mapsto m\cdot h$ and $m\mapsto m\cdot
h_x$ on $\aff\Delta$ vanish on the same supporting hyperplane, so they are proportional;
both are non-negative on $K$, which has points off that hyperplane, so the factor is
positive, and by Lemma~\ref{lem:fulldim}(a) $h$ is a positive multiple of $h_x$. Thus
$F_x=\R_+h_x$. Since $x$ has full support, $x\cdot h_x=0$ and $h_x\neq0$, some coordinate
$(h_x)_j$ is negative, and $x$ is visible from $j$ by Lemma~\ref{lem:V1}. Points of
$\relint K$ are visible by Lemma~\ref{lem:V1}. No vertex of $\Delta$ lies in $K$, so
condition (iii) of Theorem~\ref{thm:vis} holds.
\end{proof}

Examples of such $K$, whenever they lie in $\relint\Delta$, are Euclidean balls
$\{m\in\aff\Delta:|m-c|\le r\}$; outer parallel bodies $\{m\in\aff\Delta:\operatorname{dist}
(m,L)\le r\}$ of a compact convex $L\subseteq\aff\Delta$ with $r>0$ (each boundary point
lies on the boundary of a ball of radius $r$ contained in the body, and any supporting
hyperplane there supports that ball); and sublevel sets $\{m\in\aff\Delta:G(m)\le r\}$ of
a differentiable convex function $G$ on an open subset of $\aff\Delta$ with $r>\min G$
(the normal cone at a boundary point is spanned by the gradient \cite[Cor.~23.7.1]
{Rockafellar}). The last class includes ellipsoids and balls $\{m:\mathrm{KL}(m\|q)\le r\}$
for the Kullback--Leibler divergence, $q\in\relint\Delta$.

\begin{corollary}\label{cor:n2}
If $n=2$, every full-dimensional credal set $K\subseteq\Delta$ is covered by a
veto--Brier ranking satisfying the hypotheses of Theorem~\ref{thm:cover}.
\end{corollary}

\begin{proof}
$K$ is a segment with endpoints $x,y$, $x_1<y_1$. If $y\neq e_1$, the points
$y+s(y-e_1)$ have first coordinate $y_1-s(1-y_1)\in(x_1,y_1)$ for small $s>0$, so they lie in
$\relint K$ and $y$ is visible from $1$. Symmetrically $x$ is $e_2$ or visible from $2$.
Apply Corollary~\ref{cor:poly}.
\end{proof}

\begin{remark}\label{rem:qual}
Three qualifications in these statements matter.
\begin{enumerate}[label=(\arabic*)]
  \item \emph{Strict, not weak, visibility.} Visible points are weakly visible, but not
  conversely. Let $K=\conv\{v,w,x\}$ with $v=(\frac25,\frac3{10},\frac3{10})$,
  $w=(\frac15,\frac25,\frac25)$ and $x=(\frac9{20},\frac9{20},\frac1{10})$. Its facet
  normals at $v$ are $(0,1,-1)$ (through $v,w$) and $(-21,19,9)$ (through $v,x$). Both have
  non-positive first coordinate, so $v$ is weakly visible from world $1$; but
  $w=v+\frac13(v-e_1)$, so the ray from $e_1$ through $v$ runs along the edge $[v,w]$ and
  misses $\relint K$. By Lemma~\ref{lem:poly}, $v$ is not visible from worlds $2$ or $3$
  either. So $v$ is tangential from world $1$ and is not visible; by
  Lemma~\ref{lem:visbasic}(a) it lies in no set $\conv(K''\cup\{e_j\})$ with $K''$ a
  compact subset of $\relint K$. Tangential points count as not visible in
  Theorems~\ref{thm:cover} and~\ref{thm:vis}, and also in Theorem~\ref{thm:reach} below.
  \item \emph{Uniformity.} Theorem~\ref{thm:vis} requires visibility on $\cl(\ext K)$, not
  only on $\ext K$, together with condition (iii)(b) at vertices of $\Delta$. For
  polytopes both reduce to visibility of the vertices.
  \item \emph{Dimension and interiority.} Full-dimensionality of $K$ is assumed throughout;
  it is forced by (LF) on these domains (Proposition~\ref{prop:lowdim}). The assumption
  $K\subseteq\relint\Delta$ is used only in Corollary~\ref{cor:smooth} and in
  Section~\ref{sec:limits}.
\end{enumerate}
\end{remark}

\section{Examples}\label{sec:examples}

Interior $\varepsilon$-contaminations, including the vacuous model, are covered by
Corollary~\ref{cor:contam}, in every number of worlds. Smooth credal sets in
$\relint\Delta$ and all full-dimensional credal sets on two worlds are covered by
Corollaries~\ref{cor:smooth} and~\ref{cor:n2}.

\subsection{Total-variation balls}

\begin{proposition}\label{prop:tv}
Let $p\in\relint\Delta$, $0<\delta<\min_jp_j$, and $K\coloneqq\{m\in\Delta:\frac12\sum_i
|m_i-p_i|\le\delta\}$. Then $K$ is a full-dimensional polytope in $\relint\Delta$ whose
vertices are all visible, and $K$ is covered, for every $n\ge2$.
\end{proposition}

\begin{proof}
Every $m\in\aff\Delta$ with $\frac12\sum|m_i-p_i|\le\delta$ has $m_j\ge p_j-\delta>0$, so
$K=p+\delta Q$ with $Q\coloneqq\{d:\one\cdot d=0,\ \sum|d_i|\le2\}$, and $K\subseteq
\relint\Delta$ is full-dimensional. If $d\in Q$, $d\ne0$, write $d=d^+-d^-$ with
$\one\cdot d^+=\one\cdot d^-=:\sigma\le1$; then $d=\sum_{i\neq j}\frac{d^+_id^-_j}\sigma
(e_i-e_j)$, with weights summing to $\sigma$, and $0=\frac12(e_1-e_2)+\frac12(e_2-e_1)$.
So $Q=\conv\{e_i-e_j:i\neq j\}$, and every vertex of $K$ has the form $v=p+\delta(e_i-e_j)$.
For small $s>0$ the point $y\coloneqq v+s(v-e_i)$ satisfies
\[
  y-p=(1+s)\delta(e_i-e_j)+s(p-e_i),\qquad
  \tfrac12\textstyle\sum_k|y_k-p_k|=\delta+s(\delta-p_j)<\delta ,
\]
by a direct computation of the coordinates $i$, $j$ and $k\notin\{i,j\}$. So $y\in\relint K$
and $v$ is visible from $i$. Apply Corollary~\ref{cor:poly}.
\end{proof}

\subsection{Lower probabilities}

Let $\underline P:2^\Omega\to[0,1]$ with $\underline P(\emptyset)=0$ and $\underline
P(\Omega)=1$, and put $K\coloneqq\{m\in\aff\Delta:m(A)\ge\underline P(A)\text{ for all
}A\subseteq\Omega\}$, where $m(A)\coloneqq\sum_{k\in A}m_k$. The constraints on singletons
give $K\subseteq\Delta$. An event $A$ is \emph{active} at $x\in K$ if $x(A)=\underline
P(A)$.

\begin{lemma}\label{lem:events}
Suppose $K$ is full-dimensional. A point $x\in K$ with $x\neq e_j$ is visible from $j$ if
and only if every active event $A$ with $\emptyset\neq A\neq\Omega$ satisfies $j\notin A$
and $\underline P(A)>0$.
\end{lemma}

\begin{proof}
On $\aff\Delta$, $m(A)-\underline P(A)=m\cdot g_A$ with $g_A\coloneqq1_A-\underline
P(A)\one$, so Lemma~\ref{lem:poly} applies. $g_\emptyset=g_\Omega=0$, and $g_A\neq0$ for
proper nonempty $A$. For such $A$, $\underline P(A)<1$: otherwise $m(\Omega\setminus A)=0$
on $K$, and $K$ would not be full-dimensional. Hence $(g_A)_j=1-\underline P(A)>0$ if
$j\in A$, and $(g_A)_j=-\underline P(A)$ if $j\notin A$.
\end{proof}

For a permutation $\sigma$ of $\Omega$ let $S_k\coloneqq\{\sigma(1),\dots,\sigma(k)\}$ and
define the marginal vector $v^\sigma$ by $v^\sigma_{\sigma(k)}\coloneqq\underline P(S_k)
-\underline P(S_{k-1})$. If $\underline P$ is supermodular, that is, $\underline P(A\cup
B)+\underline P(A\cap B)\ge\underline P(A)+\underline P(B)$, then $K$ is the core of
$\underline P$ and its extreme points are exactly the vectors $v^\sigma$
\cite{Shapley,Ichiishi}.

\begin{proposition}\label{prop:bel}
Let $\rho$ be a mass function on the nonempty subsets of $\Omega$, $\mathrm{Bel}(A)
\coloneqq\sum_{\emptyset\neq F\subseteq A}\rho(F)$, and suppose $\rho(\Omega)>0$ and
$\rho(\{k\})>0$ for every $k$. Then the credal set $K$ of $\mathrm{Bel}$ is a
full-dimensional polytope in $\relint\Delta$, each vertex $v^\sigma$ is visible from
$\sigma(n)$, and $K$ is covered.
\end{proposition}

\begin{proof}
Belief functions are supermodular \cite{Shafer}, so $\ext K=\{v^\sigma\}$. Since
$\mathrm{Bel}(S_k)-\mathrm{Bel}(S_{k-1})$ sums $\rho(F)$ over the $F$ whose
$\sigma$-last element is $\sigma(k)$,
\[
  v^\sigma(A)=\sum_{F:\,\mathrm{last}_\sigma(F)\in A}\rho(F)\ \ge\
  \sum_{F\subseteq A}\rho(F)=\mathrm{Bel}(A).
\]
If $\sigma(n)\in A\neq\Omega$, the set $F=\Omega$ contributes to the left side but not
the right, so $v^\sigma(A)\ge\mathrm{Bel}(A)+\rho(\Omega)>\mathrm{Bel}(A)$. Hence every
active proper event at $v^\sigma$ omits $\sigma(n)$, and it has $\mathrm{Bel}(A)\ge
\sum_{k\in A}\rho(\{k\})>0$. For each proper nonempty $A$ some $\sigma$ has $\sigma(n)
\in A$, so the average of all $v^\sigma$ satisfies every constraint strictly; it lies in
$\relint K$, and $K$ is full-dimensional. Also $m_k\ge\rho(\{k\})>0$ on $K$. By
Lemma~\ref{lem:events} each $v^\sigma$ is visible from $\sigma(n)$; apply
Corollary~\ref{cor:poly}.
\end{proof}

\begin{proposition}\label{prop:2mon}
Suppose $\underline P$ is \emph{strictly supermodular}: $\underline P(A\cup B)+
\underline P(A\cap B)>\underline P(A)+\underline P(B)$ whenever neither of $A,B$ contains
the other. Suppose also $\underline P(\{k\})>0$ for every $k$. Then $K$ is a
full-dimensional polytope in $\relint\Delta$, each vertex $v^\sigma$ is visible from
$\sigma(n)$, and $K$ is covered.
\end{proposition}

\begin{proof}
$\ext K=\{v^\sigma\}$ by supermodularity. Since $\underline P\ge0$ is supermodular, it is
monotone: for $A\subseteq B$, $\underline P(B)\ge\underline P(A)+\underline P(B\setminus
A)\ge\underline P(A)$. Let $A$ be nonempty, with elements $\sigma(k_1),\dots,\sigma(k_r)$,
$k_1<\dots<k_r$, and put $C_t\coloneqq A\cap S_{k_t}$. Applying supermodularity to
$S_{k_t-1}$ and $C_t$, whose union is $S_{k_t}$ and whose intersection is $C_{t-1}$,
\[
  \underline P(S_{k_t})-\underline P(S_{k_t-1})\ \ge\ \underline P(C_t)-\underline
  P(C_{t-1}),
\]
strictly unless $S_{k_t-1}\subseteq C_t$, that is, unless $S_{k_t-1}\subseteq A$ (note
$C_t\not\subseteq S_{k_t-1}$). Summing over $t$ gives $v^\sigma(A)\ge\underline P(A)$,
with equality only if $S_{k_r-1}\subseteq A\subseteq S_{k_r}$, so only if $A=S_{k_r}$ is an
initial segment of $\sigma$. The proper initial segments omit $\sigma(n)$ and have
$\underline P(S_k)\ge\underline P(\{\sigma(1)\})>0$. Each proper nonempty $A$ is not an
initial segment for some $\sigma$ (one with $\sigma(n)\in A$), so averaging the
$v^\sigma$ gives a point of $\relint K$, and $K$ is full-dimensional; $K\subseteq
\relint\Delta$ as before. Apply Lemma~\ref{lem:events} and Corollary~\ref{cor:poly}.
\end{proof}

\begin{corollary}\label{cor:distortion}
Let $p\in\relint\Delta$ and let $f:[0,1]\to[0,1]$ be strictly convex with $f(0)=0$,
$f(1)=1$ and $f>0$ on $(0,1]$. Then the credal set of $\underline P\coloneqq f\circ p$,
$\underline P(A)=f(p(A))$, is covered. This includes $f(x)=(1-\tau)x/(1-\tau x)$,
$0<\tau<1$, the event lower probability of the constant-odds-ratio model with reference
$p$ \cite{Walley}.
\end{corollary}

\begin{proof}
Let $A,B$ be non-nested, $x\coloneqq p(A\cap B)$, $a\coloneqq p(A\setminus B)>0$ and
$b\coloneqq p(B\setminus A)>0$. With $\lambda\coloneqq b/(a+b)\in(0,1)$, $x+a=\lambda x+
(1-\lambda)(x+a+b)$ and $x+b=(1-\lambda)x+\lambda(x+a+b)$. Strict convexity gives
$f(x+a)+f(x+b)<f(x)+f(x+a+b)$, which is strict supermodularity. Also $f(p_k)>0$. Apply
Proposition~\ref{prop:2mon}. For the example, $f''(x)=2\tau(1-\tau)(1-\tau x)^{-3}>0$.
\end{proof}

The covered set in the constant-odds-ratio case is the core of the event lower
probability. It differs from the density-ratio class $\{m:m_i/p_i\le km_j/p_j\}$,
$k=1/(1-\tau)$: for $n=3$, $p$ uniform and $\tau=\frac12$, the core contains
$(\frac15,\frac3{10},\frac12)$, which violates $m_3\le2m_1$. Density-ratio classes are not
reached by this construction (Corollary~\ref{cor:out}).

\subsection{Probability intervals}

\begin{proposition}\label{prop:pri}
Let $K\coloneqq\{m\in\Delta:l_i\le m_i\le u_i\}$ be full-dimensional with $K\subseteq
\relint\Delta$. A vertex of $K$ is visible if and only if at most one of its coordinates is
at its upper bound. Hence $K$ is covered if every vertex, with at most one exception, has
at most one coordinate at its upper bound.
\end{proposition}

\begin{proof}
On $\aff\Delta$ the constraints are $m\cdot(e_i-l_i\one)\ge0$ and $m\cdot(u_i\one-e_i)
\ge0$; at a vertex $v\in\relint\Delta$ the constraints $m_k\ge0$ are inactive. Let $L$ and
$H$ be the sets of coordinates of $v$ at their lower and upper bounds. For $i\in H$,
$u_i=v_i<1$, so $u_i\one-e_i$ is negative exactly at coordinate $i$. For $i\in L$,
$l_i=v_i>0$, so $e_i-l_i\one$ is negative exactly off $i$. By Lemma~\ref{lem:poly}, $v$ is
visible from $j$ iff $H\subseteq\{j\}$ and $j\notin L$. If $|H|\ge2$ no $j$ works. If
$H=\{i\}$, then $j=i\notin L$ works, since $l_i<u_i$ by full-dimensionality. If
$H=\emptyset$, then $L\neq\Omega$ (otherwise $\sum_il_i=1$ and $K=\{l\}$), and any $j\notin
L$ works. Apply Corollary~\ref{cor:poly}.
\end{proof}

\begin{corollary}\label{cor:hexagon}
Let $n=3$. Every probability-interval model $K\subseteq\relint\Delta$ that is a hexagon is
covered by a ranking satisfying the hypotheses of Theorem~\ref{thm:cover}.
\end{corollary}

\begin{proof}
Each edge of $K$ lies on one of the six lines $m_i=l_i$, $m_i=u_i$, and the six edges have
distinct outward normals, so each line carries exactly one edge. Within $\aff\Delta$ the
outward normal of the edge on $m_i=l_i$ is $-\pi(e_i)$ and that on $m_i=u_i$ is $\pi(e_i)$,
where $\pi$ is the orthogonal projection onto $\one^\perp$. The vectors $\pi(e_i)$ have
equal length and pairwise angles of $120^\circ$, so in angular order the six normals are
$\pi(e_1),-\pi(e_3),\pi(e_2),-\pi(e_1),\pi(e_3),-\pi(e_2)$, alternating in sign. Consecutive
edges meet at a vertex, and no vertex lies on a third line (a supporting line through a
vertex that carries an edge contains that vertex in the edge, and a vertex lies on exactly
two edges). So every vertex has exactly one coordinate at its upper bound. Apply
Proposition~\ref{prop:pri}.
\end{proof}

\subsection{Near the reversed triangle}

The \emph{reversed triangle} is $R\coloneqq\{m\in\Delta_3:m_i\le\frac25\}$, with vertices
$v_1=(\frac15,\frac25,\frac25)$, $v_2=(\frac25,\frac15,\frac25)$ and
$v_3=(\frac25,\frac25,\frac15)$. It is the pari-mutuel model with uniform reference and
$\delta=\frac15$, and lies in $\relint\Delta$. No vertex of $R$ is visible
(Corollary~\ref{cor:out}). Nearby targets are covered.

\begin{proposition}\label{prop:trunc}
For $0<\eta<\frac1{10}$ let $R_\eta\coloneqq\{m\in\Delta:\frac15+\eta\le m_i\le\frac25\}$.
\begin{enumerate}[label=(\alph*)]
  \item $R_\eta$ is covered by a ranking satisfying the hypotheses of
  Theorem~\ref{thm:cover}, and $d_H(R_\eta,R)=\eta\sqrt{3/2}$.
  \item $P_\eta\coloneqq\conv(R_\eta\cup\{v_1\})$ is covered. It coincides with $R$ on
  $\{m_1\le\frac15+\eta\}$, and $d_H(P_\eta,R)=\eta\sqrt{3/2}$.
\end{enumerate}
For $\eta=\frac1{10}$, $R_\eta=\{m:m_i\ge\frac3{10}\}=\frac9{10}u+\frac1{10}\Delta$, $u$
uniform, is covered by Corollary~\ref{cor:contam}.
\end{proposition}

\begin{proof}
(a) $R_\eta\subseteq\relint\Delta$ contains $u$ in its relative interior. A vertex has two
active constraints; if two coordinates were at the upper bound $\frac25$, the third would
be $\frac15<\frac15+\eta$. So every vertex has at most one coordinate at its upper bound,
and Proposition~\ref{prop:pri} applies (with no exception). The distance from a point of
$R$ to $R_\eta$ is a convex function, so $d_H(R_\eta,R)$ is attained at a vertex of $R$,
by symmetry at $v_1$. Since $R_\eta\subseteq\{m_1\ge\frac15+\eta\}$ and $|\pi(e_1)|=
\sqrt{2/3}$, $\operatorname{dist}(v_1,R_\eta)\ge\eta\sqrt{3/2}$, with equality at
$(\frac15+\eta,\frac25-\frac\eta2,\frac25-\frac\eta2)\in R_\eta$.

(b) By (a) and Theorem~\ref{thm:vis} there are $K'\subseteq\relint R_\eta$ and
$A_i\subseteq\conv(K'\cup\{e_i\})$ with $R_\eta=\conv(K'\cup\bigcup A_i)$. Then
$P_\eta=\conv(K'\cup\bigcup A_i\cup\{v_1\})$ is full-dimensional and $K'\subseteq
\relint R_\eta\subseteq\relint P_\eta$, so Theorem~\ref{thm:exempt} applies.
$R\cap\{m_1\le\frac15+\eta\}$ is the triangle with vertices $v_1$,
$(\frac15+\eta,\frac25,\frac25-\eta)$ and $(\frac15+\eta,\frac25-\eta,\frac25)$, whose last
two vertices lie in $R_\eta$; so it is contained in $P_\eta\subseteq R$. The distance
claim follows as in (a), at $v_2$.
\end{proof}

\section{Dense approximation}\label{sec:dense}

\begin{proposition}\label{prop:dense}
Let $K\subseteq\Delta$ be a full-dimensional credal set, $u_0\in\relint K$ and
$s\in(0,1)$. Put $K'_s\coloneqq u_0+s(K-u_0)$, $A_i\coloneqq K\cap\conv(K'_s\cup\{e_i\})$
and $K_s\coloneqq\conv(K'_s\cup A_1\cup\dots\cup A_n)$. Then $K'_s\subseteq\relint K_s$,
$K_s\subseteq K$,
\[
  d_H(K_s,K)\le(1-s)\max_{k\in K}|k-u_0|,
\]
and $\VB(K'_s,A)$ covers $K_s$ on $\M_{D_{K'_s}}$. In particular every full-dimensional
credal set is a Hausdorff limit of covered polytopes.
\end{proposition}

\begin{proof}
$K'_s$ is compact, convex, full-dimensional and contained in $\relint K$, since
$(1-s)u_0+sk\in\relint K$ for $k\in K$. Each $A_i$ is compact and contained in
$\conv(K'_s\cup\{e_i\})$, and $K'_s\subseteq K_s\subseteq K$. For $k\in K$, the point
$u_0+s(k-u_0)\in K_s$ lies at distance $(1-s)|k-u_0|$ from $k$, which gives the bound.

Suppose $x\in K'_s\cap\rbd K_s$. There is a supporting hyperplane of $K_s$ at $x$, given by
Lemma~\ref{lem:fulldim}(a) as $c\cdot m\le c\cdot x$ on $K_s$ with $m\mapsto c\cdot m$
non-constant on $\aff\Delta$, that is, $c$ not a multiple of $\one$. Since $x\in\relint
K\subseteq\relint\Delta$ has full support, $c\cdot x<\max_ic_i\eqqcolon c_{i^\ast}$. For
small $t>0$, $y_t\coloneqq(1-t)x+te_{i^\ast}\in K$ because $x\in\relint K$, so $y_t\in
A_{i^\ast}\subseteq K_s$; but $c\cdot y_t=c\cdot x+t(c_{i^\ast}-c\cdot x)>c\cdot x$, a
contradiction. Hence $K'_s\subseteq\relint K_s$, $K_s$ is full-dimensional, and
Theorem~\ref{thm:cover} applies. If $K$ is a polytope, so is each $A_i$ (an intersection
of polytopes), and hence so is $K_s$. Applying this to full-dimensional polytopes
approximating a general $K$ gives the last claim.
\end{proof}

The domain $D_{K'_s}$ depends on $s$: the approximation is by targets covered on their own
domains. For the reversed triangle, Proposition~\ref{prop:trunc} gives an explicit family
with exact distances.

\section{Limits of the construction}\label{sec:limits}

This section describes what veto--Brier rankings on domains $\M_{D_{K'}}$ cannot reach.
Every statement concerns this construction only; none concerns coverage by rankings of
other forms or on other domains.

\begin{lemma}\label{lem:sublinear}
Let $K'\subseteq\Delta$ be nonempty, compact and convex, and let $\varphi:\R^n\to\R$ be
non-negative, convex and positively homogeneous, with $\varphi(h)=0$ whenever $h\in
D_{K'}$ and $h_i\ge0$. Then $\varphi(h)=\max_{v\in B}(-v\cdot h)$ for a compact convex $B$
with $0\in B\subseteq\cone(K'\cup\{e_i\})$.
\end{lemma}

\begin{proof}
A finite convex positively homogeneous function on $\R^n$ is continuous, and is the
support function of a nonempty closed convex set $\tilde B$ \cite[Thm.~13.2]{Rockafellar},
which is bounded because the function is finite: $\varphi(h)=\max_{y\in\tilde B}y\cdot h$.
Put $B\coloneqq-\tilde B$, compact and convex. If $0\notin B$, separation
gives $h$ with $-v\cdot h<0$ for all $v\in B$, so $\varphi(h)<0$; hence $0\in B$. Let
$V\coloneqq D_{K'}\cap\{h_i\ge0\}$. The hypothesis says $v\cdot h\ge0$ for $v\in B$,
$h\in V$, that is, $B\subseteq V^\ast$, the dual cone. Now $\cl V=\{h:k\cdot h\ge0\text{
for }k\in K',\ h_i\ge0\}$ (if $h$ is in the right side, $h+t\one\in V$ for $t>0$), which is
the dual cone of $Q\coloneqq\cone(K'\cup\{e_i\})$. $Q=\R_+\conv(K'\cup\{e_i\})$ is closed,
as $\conv(K'\cup\{e_i\})$ is compact and avoids $0$. By the bipolar theorem
$V^\ast=(\cl V)^\ast=Q^{\ast\ast}=Q$.
\end{proof}

So the rankings $\VB(K',B)$ are exactly the two-level rankings on $\M_{D_{K'}}$ whose
first level is a non-negative sublinear integral veto that vanishes on correct
acceptances, and whose second level is the error $b_i$.

\begin{remark}\label{rem:convexU}
Restricting to $U=D_{K'}$ loses nothing among open convex domains. Let $U\subseteq\R^n$ be
an open convex cone containing every gamble with positive coordinates, with $U\cap(-U)=
\emptyset$, and let $C\coloneqq\cl U$. Then $C\neq\R^n$, so its dual cone $C^\ast$ is
nonzero and contained in $\R^n_+$; put $K''\coloneqq C^\ast\cap\Delta$, a nonempty compact
convex set. For a closed convex cone with nonempty interior, $\inter C=\{h:y\cdot h>0\text{
for all }y\in C^\ast\setminus\{0\}\}$: if $h\in\inter C$ and $y\in C^\ast\setminus\{0\}$,
then $h-\epsilon y\in C$ for small $\epsilon>0$, so $y\cdot h\ge\epsilon|y|^2>0$;
conversely, if the minimum $c$ of $y\cdot h$ over unit vectors $y\in C^\ast$ is positive,
the ball of radius $c$ about $h$ lies in $C^{\ast\ast}=C$. Since $U$ is open and convex,
$U=\inter C=D_{K''}$.
\end{remark}

\begin{proposition}\label{prop:lowdim}
If $K\subseteq\Delta$ is a credal set that is not full-dimensional, then for no nonempty
compact convex $K'\subseteq\Delta$ does $\M_{D_{K'}}$ satisfy (LF) at $D_K$. In particular
no precise model $\{p\}$ is covered by a ranking on such a domain.
\end{proposition}

\begin{proof}
The linear span of $K$ has dimension $\dim K+1<n$, so some $h\neq0$ has $m\cdot h=0$ for
all $m\in K$. Then $\pm h\in C_K\setminus D_K=\partial D_K$ (Lemma~\ref{lem:D}(b)). If (LF)
held with an open $N$, then, as in the proof of Lemma~\ref{lem:LF}(b), $N\subseteq D_{K'}$,
so $h$ and $-h$ would both lie in $D_{K'}$, contradicting Lemma~\ref{lem:D}(a).
\end{proof}

\begin{theorem}\label{thm:reach}
Let $K\subseteq\relint\Delta$ be a full-dimensional credal set. The following are
equivalent.
\begin{enumerate}[label=(\alph*)]
  \item For some nonempty compact convex $K'\subseteq\Delta$ and veto profile $B$ over $K'$,
  $D_K$ is an undominated model of $\VB(K',B)$ on $\M_{D_{K'}}$ and (LF) holds at $D_K$.
  \item For some $K'$ and $B$, $\VB(K',B)$ covers $K$ on $\M_{D_{K'}}$.
  \item For some $p\in K$, the set $\ext K\setminus\{p\}$ is uniformly visible.
\end{enumerate}
\end{theorem}

\begin{proof}
(b)$\Rightarrow$(a) is immediate, and (c)$\Rightarrow$(b) is
Corollary~\ref{cor:exemptvis}.

(a)$\Rightarrow$(c). Let $U=D_{K'}$. Since $D_K\in\M_U$, Lemma~\ref{lem:LF}(b) gives
$K'\subseteq\relint K$. By Theorem~\ref{thm:family}(c), $D_K=D_{K(W_0,\beta,\alpha)}$ a.e.\
for an admissible triple, so $K=K(W_0,\beta,\alpha)$ by Lemma~\ref{lem:open}. Taking every
$v_\omega=0$ in the definition of $W$ shows $\nu(\beta)\in\hat W\subseteq K\subseteq
\relint\Delta$. So $\supp\beta=\Omega$, hence $W_0=\Omega$, there are no
$\alpha$-weights, and $\hat W=\{\nu(\beta)\}$. Thus
\[
  K=\conv\Bigl(K'\cup\bigcup_{\omega\in\Omega}\hat A_\omega\cup\{\nu(\beta)\}\Bigr).
\]
Then $\ext K\subseteq K'\cup\bigcup_\omega\hat A_\omega\cup\{\nu(\beta)\}$, and
$K'\cap\ext K=\emptyset$. By Lemma~\ref{lem:visbasic}(b), $K'\subseteq K_\delta$ for some
$\delta>0$, and $\hat A_\omega\subseteq\conv(K'\cup\{e_\omega\})\subseteq\conv(K_\delta
\cup\{e_\omega\})$ by Lemma~\ref{lem:veto}(d). So $\ext K\setminus\{\nu(\beta)\}$ is
uniformly visible.
\end{proof}

\begin{corollary}\label{cor:polyreach}
Let $K\subseteq\relint\Delta$ be a full-dimensional polytope. $K$ is covered by some
veto--Brier ranking on some domain $\M_{D_{K'}}$ if and only if at most one vertex of $K$
is not visible. Tangential vertices count as not visible.
\end{corollary}

\begin{proof}
For a finite set of points of $K$, uniform visibility is equivalent to visibility of each
point (Lemma~\ref{lem:visbasic}(a), and the proof of Corollary~\ref{cor:poly}); and no
vertex of $\Delta$ lies in $K$. Apply Theorem~\ref{thm:reach}.
\end{proof}

\begin{corollary}\label{cor:out}
Let $n\ge3$. None of the following credal sets is covered by a veto--Brier ranking on any
domain $\M_{D_{K'}}$.
\begin{enumerate}[label=(\alph*)]
  \item \emph{Pari-mutuel models} $K=\{m\in\Delta:m_i\le(1+\delta)p_i\}$ with
  $p\in\relint\Delta$, $\delta>0$ and $K\subseteq\relint\Delta$; in particular the reversed
  triangle $R$. No vertex is visible.
  \item \emph{Density-ratio classes} $K=\{m\in\Delta:m_i/p_i\le k\,m_j/p_j\text{ for all
  }i,j\}$ with $p\in\relint\Delta$ and $k>1$. These include the credal sets of
  constant-odds-ratio models \cite{Walley}. At least three vertices are not visible.
  \item Probability-interval models $K\subseteq\relint\Delta$ with at least two vertices
  having two or more coordinates at their upper bounds.
\end{enumerate}
\end{corollary}

\begin{proof}
In each case $K$ is a full-dimensional polytope in $\relint\Delta$; we apply
Corollary~\ref{cor:polyreach} and Lemma~\ref{lem:poly}. We use that a point $v$ of
$\{m:\one\cdot m=1,\ g_f\cdot m\ge0\}$ is a vertex iff $\one$ and the active $g_f$ span
$\R^n$, equivalently iff $v$ is the only point of $\aff\Delta$ at which all active
constraints are equalities \cite[Sec.~8.5]{Schrijver}.

(a) Put $u_i\coloneqq(1+\delta)p_i$. $p$ satisfies every constraint strictly, so $K$ is
full-dimensional. At a vertex $v$ the constraints $m_k\ge0$ are inactive, and with $H$ the
set of $i$ with $v_i=u_i$, the span of $\one$ and the normals $u_i\one-e_i$ ($i\in H$)
is the span of $\one$ and $e_i$ ($i\in H$); so $|H|\ge n-1\ge2$. Each $u_i\one-e_i$ with
$i\in H$ is negative exactly at $i$ (as $u_i=v_i<1$), so by Lemma~\ref{lem:poly} $v$ is
not visible. A full-dimensional polytope in $\aff\Delta$ has at least $n\ge3$ vertices.

(b) If $m_j=0$ for some $m\in K$ then $m=0$; so $K\subseteq\relint\Delta$, and $p$
satisfies all constraints strictly. Write the constraints as $m\cdot g^{ij}\ge0$ with
$g^{ij}\coloneqq(k/p_j)e_j-(1/p_i)e_i$, negative exactly at $i$. For $\emptyset\neq H
\subsetneq\Omega$ let $w^H\coloneqq\nu\bigl(\sum_{i\in H}kp_ie_i+\sum_{i\notin H}p_ie_i
\bigr)\in K$. The active constraints at $w^H$ are those with $i\in H$, $j\notin H$. If all
of them hold with equality at $m\in\aff\Delta$, the ratios $m_i/p_i$ take a common value
$c$ on $H$ and $c/k$ off $H$, so $m=w^H$; hence $w^H$ is a vertex. If $|H|\ge2$, two
active normals are negative at different coordinates, so $w^H$ is not visible. Distinct
$H$ give distinct vertices (the coordinates of maximal ratio form $H$), and there are
$2^n-n-2\ge3$ sets $H$ with $2\le|H|\le n-1$.

(c) This is Proposition~\ref{prop:pri}.
\end{proof}

\begin{proposition}\label{prop:contamzero}
Let $n\ge3$, $0<\varepsilon<1$, and $p\in\Delta\setminus\relint\Delta$. Then
$K\coloneqq(1-\varepsilon)p+\varepsilon\Delta$ is not covered by a veto--Brier ranking on
any domain $\M_{D_{K'}}$. More precisely, there are no $K'$ and $B$ as in
Theorem~\ref{thm:reach}(a).
\end{proposition}

\begin{proof}
$K$ is a full-dimensional simplex with vertices $v_i=(1-\varepsilon)p+\varepsilon e_i$.
Suppose $K'$ and $B$ are as in Theorem~\ref{thm:reach}(a). As in that proof, $K'\subseteq
\relint K\subseteq\relint\Delta$ and $K=K(W_0,\beta,\alpha)=\conv T$ with $T\coloneqq
K'\cup\bigcup_{\omega\in W_0}\hat A_\omega\cup\hat W$, for an admissible triple. Points of
$K'$ have full support. Every point of $\cone(K'\cup\{e_\omega\})$ is $0$, has full
support, or is a positive multiple of $e_\omega$; so every point of $\hat A_\omega$ has
full support or equals $e_\omega$. Let $S\coloneqq\supp p$ and $Z\coloneqq\Omega\setminus
S\neq\emptyset$.

\emph{Case $|S|\ge2$.} Then $e_\omega\notin K$ for every $\omega$ (as $e_\omega-(1-
\varepsilon)p$ has a negative coordinate). The face $G\coloneqq\{m\in K:m_k=0\text{ for }
k\in Z\}=\conv\{v_i:i\in S\}$ has $|S|\ge2$ extreme points, each with support exactly $S$.
Since $m\mapsto\sum_{k\in Z}m_k$ is non-negative on $K$ and vanishes on $G$, every
representation of a point of $G$ as a convex combination of points of $T$ uses only points
of $T\cap G$; so $G=\conv(T\cap G)$, and by the above $T\cap G=\hat W\cap G$. A point $\nu(w)$, $w=\beta+\sum_{\omega\notin W_0}\alpha_\omega v_\omega$, lies in $G$
iff $\supp w\subseteq S$, which forces $\supp\beta\subseteq S$ and, for each $\omega$ with
$\alpha_\omega>0$, either $v_\omega=0$ or $v_\omega=t_\omega e_\omega$ with $\omega\in S$.
With $T_\omega\coloneqq\max\{t:te_\omega\in B_\omega\}$, the set of such $w$ is the box
$\beta+\sum_{\omega\in S\setminus W_0}\alpha_\omega[0,T_\omega]e_\omega$. Since $\nu$ maps
convex combinations to convex combinations, $G=\conv\nu(\text{corners})$, so $\ext G$
consists of normalised corners $\beta+\sum_{\omega\in M}\alpha_\omega T_\omega e_\omega$,
$M\subseteq S\setminus W_0$. A corner has support $\supp\beta\cup\{\omega\in M:
\alpha_\omega T_\omega>0\}$, and at most one corner point has support exactly $S$. This
contradicts $|\ext G|\ge2$.

\emph{Case $|S|=1$}, say $p=e_1$. For $i\ge2$, $v_i=(1-\varepsilon)e_1+\varepsilon e_i$ is
extreme with support $\{1,i\}$, so it is neither in $K'$ nor in any $\hat A_\omega$; hence
$v_i\in\hat W$. Write $v_i=\nu(w^i)$. Since $\supp w^i=\{1,i\}$ for $i=2,3$, $\supp\beta
\subseteq\{1\}$, so $\beta=\beta_1e_1$ with $\beta_1>0$ and $1\in W_0$. Each $w^i$ is then
$\beta_1e_1+\alpha_it_ie_i$ with $i\notin W_0$, $t_ie_i\in B_i$, all other terms zero; and
$\nu(w^i)=v_i$ forces $\alpha_it_i=a\beta_1$ with $a\coloneqq\varepsilon/(1-\varepsilon)$.
Choosing $v_2=t_2e_2$, $v_3=t_3e_3$ and all other $v_\omega=0$ gives $w\coloneqq\beta_1(e_1
+ae_2+ae_3)\in W$. But the first coordinate of $\nu(w)$ is $1/(1+2a)=(1-\varepsilon)/
(1+\varepsilon)<1-\varepsilon\le\min_{m\in K}m_1$, so $\nu(w)\notin K$, contradicting
$\hat W\subseteq K$.
\end{proof}

\begin{remark}\label{rem:scope}
Corollary~\ref{cor:out} and Propositions~\ref{prop:lowdim} and~\ref{prop:contamzero}
describe the reach of rankings $\VB(K',B)$ on domains $\M_{D_{K'}}$, that is, of two-level
rankings with a sublinear integral veto followed by the Brier-type error, on open convex
one-sided domains. They are not claims that the credal sets in question cannot be covered
by rankings of other forms, or on other domains. For the reversed triangle, the nearest
results available from this construction are those of Proposition~\ref{prop:trunc}: every
truncation $R_\eta$ is covered, as is $R_\eta$ with one exact corner restored; and by
Theorem~\ref{thm:reach}, a target in $\relint\Delta$ reached by this construction has at
most one extreme point that is not uniformly visible.
\end{remark}

\begin{thebibliography}{99}

\bibitem{Ichiishi}
T.~Ichiishi.
\newblock Super-modularity: applications to convex games and to the greedy algorithm
for {LP}.
\newblock \emph{Journal of Economic Theory}, 25:283--286, 1981.

\bibitem{Lindenstrauss}
J.~Lindenstrauss.
\newblock A short proof of {L}iapounoff's convexity theorem.
\newblock \emph{Journal of Mathematics and Mechanics}, 15:971--972, 1966.

\bibitem{Lyapunov}
A.~Lyapunov.
\newblock Sur les fonctions-vecteurs compl\`etement additives.
\newblock \emph{Izvestiya Akademii Nauk SSSR, Seriya Matematicheskaya}, 4:465--478,
1940.

\bibitem{Rockafellar}
R.~T. Rockafellar.
\newblock \emph{Convex Analysis}.
\newblock Princeton University Press, 1970.

\bibitem{Schrijver}
A.~Schrijver.
\newblock \emph{Theory of Linear and Integer Programming}.
\newblock Wiley, 1986.

\bibitem{Shafer}
G.~Shafer.
\newblock \emph{A Mathematical Theory of Evidence}.
\newblock Princeton University Press, 1976.

\bibitem{Shapley}
L.~S. Shapley.
\newblock Cores of convex games.
\newblock \emph{International Journal of Game Theory}, 1:11--26, 1971.

\bibitem{Walley}
P.~Walley.
\newblock \emph{Statistical Reasoning with Imprecise Probabilities}.
\newblock Chapman and Hall, 1991.

\end{thebibliography}

\end{document}
